% Mechanically extracted after \mainmatter from the frozen standalone master.

\markboth{Exposé XIV. Regular Elements: Continuation}
  {Exposé XIV. Regular Elements: Continuation}
% BEGIN inputs/XIV/tex/components/00_expose_XIV_section1_opening_en.tex
% Authority: Expo14.pdf, local pp. 1--3 / combined pp. 835--837.
% Continuous English currently runs through the statement of Lemma 2.1.

\SGARefTarget{sga3:xiv:expose}

\setcounter{footnote}{0}
\footnotetext{Version xy of 1 December 2008.}

\section{Construction of Cartan subgroups and maximal tori for a
smooth algebraic group}
\SGARefTarget{sga3:xiv:section:1}

\paragraph{\textbf{Theorem 1.1.}}
\SGARefTarget{sga3:xiv:1.1}
\emph{Let \(G\) be a smooth algebraic group over a field \(k\).
Then \(G\) admits a maximal torus \(T\), and hence a Cartan subgroup
\(C=\Centr_G(T)\).}

By \SGARefLink{sga3:xii:3.2}{XII~3.2}, finding a maximal torus \(T\) of \(G\) is equivalent to
finding a Cartan subgroup \(C\) of \(G\). Moreover, since the maximal
tori of \(G\) are those of \(G^0\), we may suppose that \(G\) is
connected. We distinguish two cases.

\medskip
\noindent\textup{1\textdegree)} \emph{The field \(k\) is finite.}
Let \(\mathcal T\) be the scheme of maximal tori of \(G\)
(\SGARefLink{sga3:xii:1.10}{XII~1.10}); it is smooth over \(k\). The group \(G\) acts on
\(\mathcal T\) through inner automorphisms, and the conjugacy theorem
(\SGARefLink{sga3:xii:6.6:item:a}{XII~6.6(a)}) says that any two \(\bar k\)-rational points of
\(\mathcal T_{\bar k}\) are conjugate under
\(G_{\bar k}(\bar k)\). Since \(\mathcal T\) is smooth over \(k\),
\(\mathcal T_{\bar k}\) is smooth over \(\bar k\); it follows that
\(\mathcal T_{\bar k}\) is isomorphic to
\[
  G_{\bar k}/N,
\]
where \(N\) is the stabilizer of an element of
\(\mathcal T_{\bar k}(\bar k)\), that is, the normalizer of a
maximal torus \(\bar T\) of \(G_{\bar k}\). Consequently
\(\mathcal T\) is a homogeneous space under the action of \(G\).
A well-known theorem of Lang
(\emph{Amer.\ J.\ Math.} \textbf{78} (1956), 555--563) says that
every homogeneous space under a smooth connected algebraic group
over a finite field \(k\) has a rational point. In particular,
\(\mathcal T\) has a rational point; equivalently, \(G\) has a
maximal torus \(T\).

\medskip
\noindent\textup{2\textdegree)} \emph{The field \(k\) is infinite.}
We shall use the following lemma.

\paragraph{\textbf{Lemma 1.2.}}
\SGARefTarget{sga3:xiv:1.2}
\emph{Let \(G\) be a smooth algebraic group over a field \(k\).
Then \(G\) admits a subgroup of type \textup{(C)}
\textup{(\SGARefLink{sga3:xiii:definition:6:2}{XIII~6.2})}.}

By \SGARefLink{sga3:xiii:corollary:6:3}{XIII~6.3}, this is equivalent to saying that \(\mathfrak g\)
contains a Cartan subalgebra \(\mathfrak d\). When \(k\) is infinite
this is immediate: \(\mathfrak g\) contains a regular element \(a\),
and one takes
\[
  \mathfrak d=\Nil(a,\mathfrak g).
\]
The case of finite \(k\) is treated exactly as in the proof of
\textup{1\textdegree)} above, but it requires first constructing the
scheme \(\mathcal D\) of Cartan subalgebras of \(\mathfrak g\) and
proving that this scheme is smooth over \(k\); this will be done
below in \SGARefLink{sga3:xiv:2.16}{2.16}.

To prove \SGARefLink{sga3:xiv:1.1}{1.1} in case \textup{2\textdegree)}, which is the case now
under consideration, it is therefore enough to know \SGARefLink{sga3:xiv:1.2}{1.2} when \(k\)
is infinite. We also record the following result for later use.

\paragraph{\textbf{Lemma 1.3.}}
\SGARefTarget{sga3:xiv:1.3}
\emph{Let \(G\) be a smooth connected affine algebraic group whose
reductive center \textup{(\SGARefLink{sga3:xii:4.1}{XII~4.1} and \SGARefLink{sga3:xiv:4.4}{4.4})} is reduced to the identity
group. Then \(G\) is nilpotent if and only if its Lie algebra
\(\mathfrak g\) is nilpotent.}

This is contained in \SGARefLink{sga3:xii:4.9}{XII~4.9}.

We can now give a procedure for constructing Cartan subgroups of
\(G\). The same procedure applies when \(k\) is finite, provided
\SGARefLink{sga3:xiv:1.2}{Lemma~1.2} is granted in that case. Suppose first that \(G\) is
affine. We argue by induction on \(n=\dim G\), the assertion being
trivial when \(n=0\). Suppose therefore that \(n>0\), and that the
assertion has been proved in dimensions \(n'<n\). Let \(Z\) be the
reductive center of \(G\), and let
\[
  u\colon G\longrightarrow G'=G/Z
\]
be the canonical homomorphism. By \SGARefLink{sga3:xii:theorem:4-7:item:c}{XII~4.7(c)}, the rule
\[
  C'\longmapsto C=u^{-1}(C')
\]
is a bijection from the Cartan subgroups of \(G'\) to those of
\(G\). Replacing \(G\) by \(G'\), we may thus suppose that the
reductive center \(Z\) of \(G\) is reduced to the identity element
\textup{(the same is true for the reductive center of \(G'\), by
\SGARefLink{sga3:xii:theorem:4-7:item:b}{XII~4.7(b)})}.

By \SGARefLink{sga3:xiv:1.2}{Lemma~1.2},\footnote{The source prints ``\SGARefLink{sga3:xiv:1.3}{1.3}'' here; the cited
existence result is \SGARefLink{sga3:xiv:1.2}{Lemma~1.2}.} \(G\) admits a subgroup \(D\) of
type \textup{(C)}. Over the algebraic closure \(\bar k\) of \(k\),
the group \(D_{\bar k}\) contains a Cartan subgroup of
\(G_{\bar k}\) \textup{(\SGARefLink{sga3:xiii:remarks:6:6:part:b}{XIII~6.6(b)})}; hence every Cartan subgroup
of \(D\) is a Cartan subgroup of \(G\) \textup{(\SGARefLink{sga3:xiii:corollary:2:8:item:a}{XIII~2.8(a)})}.
We are therefore reduced to finding a Cartan subgroup of \(D\).

If \(\dim D=\dim G\), that is, if \(D=G\), then the Lie algebra of
\(G\) is a Cartan subalgebra of itself and is therefore nilpotent.
\SGARefLink{sga3:xiv:1.3}{Lemma~1.3} then shows that \(G\) is nilpotent, so \(G\) is a Cartan
subgroup of itself. If \(\dim D<\dim G\), the induction hypothesis
gives a Cartan subgroup of \(D\), which is consequently a Cartan
subgroup of \(G\). This proves the assertion when \(G\) is affine.

In the general case, let \(Z\) be the center of \(G\). Then
\(G/Z=G'\) is affine \textup{(\SGARefLink{sga3:xii:6.1}{XII~6.1})}, and the inverse image in
\(G\) of every Cartan subgroup \(C'\) of \(G'\) is a Cartan
subgroup of \(G\) \textup{(\SGARefLink{sga3:xii:6.6:item:e}{XII~6.6(e)})}. We are thus reduced to
finding a Cartan subgroup of the smooth affine algebraic group
\(G'\), which is the case already treated.

\paragraph{\textbf{Corollary 1.4.}}
\SGARefTarget{sga3:xiv:1.4}
\emph{Let \(G\) be a smooth group scheme of finite type over an
Artinian scheme \(S\). Then \(G\) admits a maximal torus \(T\), and
hence a Cartan subgroup \(C=\Centr_G(T)\). Every torus in \(G\) is
contained in a maximal torus.}

Indeed, we may suppose that \(S\) is local, with residue field
\(k\). By \SGARefLink{sga3:xiv:1.1}{1.1}, the group
\[
  G_0=G\times_S\operatorname{Spec}(k)
\]
admits a maximal torus \(T_0\). By \SGARefLink{sga3:ix:theorem:3:6}{IX~3.6}~bis and \SGARefLink{sga3:x:corollary:2:3}{X~2.3}, \(T_0\)
comes from a torus \(T\) of \(G\), which is plainly maximal. The
last assertion follows by applying the preceding result to the
centralizer of the given torus, which is smooth over \(k\)
\textup{(\SGARefLink{sga3:xi:corollary:2:4}{XI~2.4})}.

\paragraph{\textbf{Remarks 1.5.}}
\SGARefTarget{sga3:xiv:1.5}
\textup{(a)} We shall give below, in \SGARefLink{sga3:xiv:3.20}{3.20}, \SGARefLink{sga3:xiv:3.21}{3.21}, and Exposé XV,
variants of \SGARefLink{sga3:xiv:1.4}{1.4} in which \(S\) is not assumed Artinian.

\medskip
\noindent\textup{\SGAInlineRefTarget{sga3:xiv:remark:1-5:item:b}(b)} I do not know whether every algebraic group
\(G\), not necessarily smooth, over a field \(k\) admits a maximal
torus. The question arises only in characteristic \(p>0\). Applying
\SGARefLink{sga3:xiv:1.1}{1.1} to a smooth quotient group of the form \(G'=G/I\), where \(I\)
is a suitable radicial subgroup of \(G\) \textup{(for example, the
kernel of a sufficiently high power of the Frobenius homomorphism)},
and then taking in \(G\) the inverse image of a maximal torus of
\(G'\), reduces the question to the case in which
\((G_{\bar k})_{\mathrm{red}}\) is a torus; here, as always,
\(\bar k\) denotes the algebraic closure of \(k\). The answer is
easily seen to be affirmative when \(G\) is commutative, or more
generally nilpotent: in that case \(G\) has a unique maximal torus,
which may for example be constructed by descent from the maximal
torus of \(G_{\bar k}\).\begingroup
\renewcommand{\thefootnote}{*}\footnote{M. Raynaud gave a negative
answer to the question raised here; see XVII, \SGARefLink{sga3:xvii:subsection:5-9:item:c}{Example~5.9(c)}.}
\endgroup

\medskip
\noindent\textup{(c)} When \(G\) is affine and either \(k\) is
perfect or \(G\) is solvable, \SGARefLink{sga3:xiv:1.1}{1.1} was already known, by a theorem of
Rosenlicht; his proof is quite different from the one given here.

\medskip
\noindent\textup{(d)} When \(k\) is infinite, \SGARefLink{sga3:xiv:1.1}{1.1} follows from the
much more precise result, proved below in \SGARefLink{sga3:xiv:6.1}{6.1}, that the scheme
\(\mathcal T\) of maximal tori of \(G\) is a rational variety. The
method essentially combines the proof of \SGARefLink{sga3:xiv:1.1}{1.1} with an explicit
description of the structure of the scheme \(\mathcal D\) of Cartan
subalgebras of \(\mathfrak g\). To obtain the desired result, we
must first extend to an arbitrary base prescheme certain results of
XIII~4--6 \textup{(this is the purpose of the next two sections)},
and then refine the preceding construction proving \SGARefLink{sga3:xiv:1.1}{1.1}, using the
fact that every Cartan subgroup of \(G\) is contained in one and
only one subgroup of type \textup{(C)} in \(G\) \textup{(Nos.~4
and 5)}.

\section{Lie algebras over an arbitrary prescheme: regular sections
and Cartan subalgebras}
\SGARefTarget{sga3:xiv:section:2}

\paragraph{\textbf{Lemma 2.1.}}
\SGARefTarget{sga3:xiv:2.1}
\emph{Let \(A\) be a ring, let \(\mathfrak d\) be a Lie algebra over
\(A\), and for each \(s\in\operatorname{Spec}(A)\) let
\[
  \mathfrak d(s)=\mathfrak d\otimes_A k(s)
\]
be the corresponding Lie algebra over the residue field \(k(s)\).
Suppose that the \(A\)-module \(\mathfrak d\) is of finite
presentation. The following conditions are equivalent:
\begin{enumerate}
\item[\textup{(i)}] for every \(s\in\operatorname{Spec}(A)\), the
  Lie algebra \(\mathfrak d(s)\) is nilpotent;
\item[\textup{(ii)}] for every \(x\in\mathfrak d\), the endomorphism
  \(\operatorname{ad}(x)\) is nilpotent;
\item[\textup{(iii)}] there is an integer \(N\geq 0\) such that,
  for every sequence \(x_1,\ldots,x_N\) of elements of
  \(\mathfrak d\), one has
  \[
    \operatorname{ad}(x_1)\operatorname{ad}(x_2)\cdots
    \operatorname{ad}(x_N)=0.
  \]
\end{enumerate}}

% END inputs/XIV/tex/components/00_expose_XIV_section1_opening_en.tex

% BEGIN inputs/XIV/tex/components/01_expose_XIV_lemma21_proof_through_definition25_en.tex
% Authority: Expo14.pdf, local pp. 3--4 / combined pp. 837--838.
% Continuous English ends immediately before Proposition 2.6.

When \(A\) is a field, the equivalence of \textup{(i)} and
\textup{(iii)} is the definition of nilpotence, while the equivalence
of \textup{(ii)} and \textup{(iii)} is the familiar consequence of
Engel's theorem \textup{(Bourbaki, \emph{Groupes et algèbres de Lie},
Chapter I, \S4, No.~2)}. In the general case,
\textup{(iii)} plainly implies \textup{(ii)}, and \textup{(ii)}
implies \textup{(i)} by the preceding result and the fact that
\textup{(ii)} is preserved under passage to quotients and under
localization. It remains to prove that \textup{(i)} implies
\textup{(iii)}.

Suppose first that \(A\) is local Artinian, with maximal ideal
\(\mathfrak m\). Choose \(n>0\) such that \(\mathfrak m^n=0\), and
choose \(N\) such that condition \textup{(iii)} holds for
\[
  \mathfrak d(s)=\mathfrak d\otimes_A A/\mathfrak m.
\]
Taking \(N'=nN\), one sees at once that \(N'\) satisfies
\textup{(iii)} for \(\mathfrak d\).

If \(A\) is Noetherian, there are finitely many points
\(s_i\in\operatorname{Spec}(A)\) and ideals of definition
\(\mathfrak q_i\subset A_{s_i}\) such that the natural map
\[
  \mathfrak d\longrightarrow
  \prod_i\bigl(\mathfrak d\otimes_A A_i\bigr),
  \qquad A_i=A_{s_i}/\mathfrak q_i,
\]
is injective. By the preceding case, for each \(i\) there is an
integer \(N_i\) satisfying \textup{(iii)} for the Lie algebra
\(\mathfrak d\otimes_A A_i\). Taking \(N\) to be the largest of the
\(N_i\), condition \textup{(iii)} follows for \(\mathfrak d\).
Finally, the general case reduces to the Noetherian case by the
usual procedure described in EGA~IV~8.

\paragraph{\textbf{Definition 2.2.}}
\SGARefTarget{sga3:xiv:2.2}
Let \(S\) be a prescheme and let \(\mathfrak d\) be a
quasi-coherent Lie algebra over \(S\) whose underlying module is of
finite presentation. We say that \(\mathfrak d\) is
\emph{locally nilpotent} if, for every \(s\in S\), the Lie algebra
\(\mathfrak d(s)\) over \(k(s)\) is nilpotent. We say that
\(\mathfrak d\) is \emph{strictly locally nilpotent} if it is
locally free and, on every open subset \(U\) of \(S\) over which it
has constant rank \(r\), its Killing polynomial reduces to
\[
  P_{\mathfrak d}(t)=t^r.
\]

Of course, when \(\mathfrak d\) is a Lie algebra over \(S\) that is
locally free as a module, its Killing polynomial is defined as a
polynomial
\[
  P_{\mathfrak d}\in A[t],
  \qquad
  A=\Gamma\!\left(\operatorname{Sym}_{\mathcal O_S}(\mathfrak d)\right)
   \simeq \Gamma\!\left(W(\mathfrak d)\right),
\]
where \(A\) is the ring of sections of the structure sheaf of the
vector bundle \(W(\mathfrak d)\) defined by \(\mathfrak d\).

The two notions introduced in \SGARefLink{sga3:xiv:2.2}{2.2} are plainly stable under base
change and local for the fpqc topology. We record:

\paragraph{\textbf{Proposition 2.3.}}
\SGARefTarget{sga3:xiv:2.3}
\emph{If \(\mathfrak d\) is strictly locally nilpotent, then it is
locally nilpotent. The converse holds when \(\mathfrak d\) is
locally free and \(S\) is reduced.}

The proof is immediate.

\paragraph{\textbf{Definition 2.4.}}
\SGARefTarget{sga3:xiv:2.4}
Let \(S\) be a prescheme, let \(\mathfrak g\) be a Lie algebra over
\(S\) whose underlying module is locally free of finite type, and
let \(\mathfrak d\) be a Lie subalgebra of \(\mathfrak g\). We call
\(\mathfrak d\) a \emph{Cartan subalgebra} if it satisfies the
following conditions:
\begin{enumerate}
\item[\textup{\SGAInlineRefTarget{sga3:xiv:2.4:item:i}(i)}] \(\mathfrak d\) is locally a direct-factor
  submodule \textup{(and hence is itself locally free of finite
  type)};
\item[\textup{(ii)}] for every \(s\in S\), the Lie algebra
  \(\mathfrak d(s)\) is a Cartan subalgebra of
  \(\mathfrak g(s)\).
\end{enumerate}

\paragraph{\textbf{Definition 2.5.}}
\SGARefTarget{sga3:xiv:2.5}
Let \(S\) and \(\mathfrak g\) be as in \SGARefLink{sga3:xiv:2.4}{2.4}. A section \(a\) of
\(\mathfrak g\) is called \emph{quasi-regular} if, for every
\(s\in S\), the element \(a(s)\in\mathfrak g(s)\) is regular in the
Lie algebra \(\mathfrak g(s)\) over \(k(s)\). We call \(a\) a
\emph{regular section} if it is quasi-regular and is contained in a
Cartan subalgebra of \(\mathfrak g\).

The notions in \SGARefLink{sga3:xiv:2.4}{2.4} and \SGARefLink{sga3:xiv:2.5}{2.5} are again stable under base change and
local for the fpqc topology. Only the last assertion, and only for
the notion of a regular section, requires proof. It follows from the
fact that the Cartan subalgebra containing a given regular section
is uniquely determined. More precisely:

% END inputs/XIV/tex/components/01_expose_XIV_lemma21_proof_through_definition25_en.tex

% BEGIN inputs/XIV/tex/components/02_expose_XIV_proposition26_corollary27_en.tex
% Authority: Expo14.pdf, local p. 5 / combined p. 839.
% Continuous English ends immediately before Corollary 2.8.

\paragraph{\textbf{Proposition 2.6.}}
\SGARefTarget{sga3:xiv:2.6}
\emph{Let \(S\) and \(\mathfrak g\) be as in \SGARefLink{sga3:xiv:2.4}{2.4}, and let \(a\) be a
quasi-regular section of \(\mathfrak g\). There is at most one Cartan
subalgebra of \(\mathfrak g\) containing \(a\). Such a Cartan
subalgebra exists---that is, \(a\) is a regular section---if and only
if \(a\) satisfies the following condition:
\[
  \mathfrak d=\Nil(a,\mathfrak g)
\]
is locally a direct-factor submodule of \(\mathfrak g\), and
\(\operatorname{ad}(a)\) induces an automorphism of
\(\mathfrak g/\mathfrak d\). In that case \(\mathfrak d\) is the
unique Cartan subalgebra of \(\mathfrak g\) containing \(a\).}

Suppose first that \(a\) is contained in a Cartan subalgebra
\(\mathfrak d\) of \(\mathfrak g\). Then
\(\operatorname{ad}(a)_{\mathfrak g/\mathfrak d}\) is bijective on
every fiber. Since \(\mathfrak g/\mathfrak d\) is locally free of
finite type, it follows that this endomorphism is an automorphism of
\(\mathfrak g/\mathfrak d\). On the other hand, by \SGARefLink{sga3:xiv:2.1}{2.1} the
endomorphism \(\operatorname{ad}(a)_{\mathfrak d}\) is locally
nilpotent. Hence
\[
  \mathfrak d=\Nil(a,\mathfrak g),
\]
which proves both the uniqueness of \(\mathfrak d\) and the
necessity of the stated regularity criterion.

For sufficiency, observe that the hypothesis on \(a\) implies that
the formation of \(\Nil(a,\mathfrak g)\) commutes with base
extension, and in particular with passage to fibers. Thus the fibers
\(\mathfrak d(s)\) of
\(\mathfrak d=\Nil(a,\mathfrak g)\) are Cartan subalgebras of the
\(\mathfrak g(s)\). Moreover, \(\mathfrak d\) is a Lie subalgebra of
\(\mathfrak g\) by \SGARefLink{sga3:xiii:proposition:4:1}{XIII~4.1}. It is therefore a Cartan subalgebra.

\paragraph{\textbf{Corollary 2.7.}}
\SGARefTarget{sga3:xiv:2.7}
\emph{Under the hypotheses of \SGARefLink{sga3:xiv:2.4}{2.4}, let \(\mathfrak d\) be a Cartan
subalgebra of \(\mathfrak g\), let \(a\) be a section of
\(\mathfrak d\) that is regular in \(\mathfrak g\), and let \(u\) be
an automorphism of \(\mathfrak g\). Then \(u\) preserves
\(\mathfrak d\) if and only if \(u(a)\) is a section of
\(\mathfrak d\).}

Indeed, by transport of structure \(u(a)\) is then a regular section
of \(\mathfrak g\), contained in both Cartan subalgebras
\(\mathfrak d\) and \(u(\mathfrak d)\); \SGARefLink{sga3:xiv:2.6}{Proposition~2.6} therefore
shows that these two subalgebras are equal.

% END inputs/XIV/tex/components/02_expose_XIV_proposition26_corollary27_en.tex

% BEGIN inputs/XIV/tex/components/03_expose_XIV_corollary28_and_nilpotent_rank_en.tex
% Authority: Expo14.pdf, local p. 5 / combined p. 839.
% Continuous English ends immediately before Proposition 2.9.

\paragraph{\textbf{Corollary 2.8.}}
\SGARefTarget{sga3:xiv:2.8}
\emph{Under the hypotheses of \SGARefLink{sga3:xiv:2.4}{2.4}, let \(\mathfrak d\) be a Cartan
subalgebra of \(\mathfrak g\). For every \(s\in S\) such that
\(k(s)\) is infinite, there exist an open neighborhood \(V\) of
\(s\) and a regular section \(a\) of \(\mathfrak g\) over \(V\)
such that
\[
  \mathfrak d|_V=\Nil\!\left(a,\mathfrak g|_V\right);
\]
equivalently, \(a\) is a section of \(\mathfrak d|_V\).}

Indeed, since \(k(s)\) is infinite, there is a regular element
\(a_0\) of \(\mathfrak g(s)\) contained in \(\mathfrak d(s)\).
After replacing \(S\) by an open neighborhood \(U\) of \(s\), we
may extend \(a_0\) to a section \(a\) of \(\mathfrak d\) over
\(U\). The endomorphism induced by \(\operatorname{ad}(a)\) on the
fiber
\[
  \mathfrak g(s)/\mathfrak d(s)
\]
is an automorphism. Shrinking \(U\) to a smaller open neighborhood
\(V\), we may therefore suppose that
\(\operatorname{ad}(a)_{\mathfrak g/\mathfrak d}\) is an
automorphism. It follows that \(a\) is a quasi-regular section of
\(\mathfrak g|_V\) satisfying the required condition.

Let \(\mathfrak g\) be as in \SGARefLink{sga3:xiv:2.4}{2.4}. Inspection of its Killing
polynomial immediately shows that the function
\[
  s\longmapsto\text{the nilpotent rank of }\mathfrak g(s)
\]
on \(S\) is upper semicontinuous. We shall be especially interested
in the case in which this function is continuous, that is, locally
constant. Here are several variants of this property.

% END inputs/XIV/tex/components/03_expose_XIV_corollary28_and_nilpotent_rank_en.tex

% BEGIN inputs/XIV/tex/components/04_expose_XIV_proposition29_through_corollary211_en.tex
% Authority: Expo14.pdf, local pp. 5--6 / combined pp. 839--840.
% Continuous English ends immediately before Remark 2.12.

\paragraph{\textbf{Proposition 2.9.}}
\SGARefTarget{sga3:xiv:2.9}
{\itshape Let \(S,\mathfrak g\) be as in \SGARefLink{sga3:xiv:2.4}{2.4}. Consider the following
conditions:

\smallskip
\noindent
\textup{\((\mathrm C_0)\)}
The nilpotent rank of \(\mathfrak g(s)\), for \(s\in S\), is a
locally constant function of \(s\).

\smallskip
\noindent
\textup{\((\mathrm C_1)\)}
Locally for the fpqc topology, there exists a Cartan subalgebra of
\(\mathfrak g\).

\smallskip
\noindent
\textup{\((\mathrm C'_1)\)}
The same as \((\mathrm C_1)\), with ``fpqc topology'' replaced by
``étale topology''.

\smallskip
\noindent
\textup{\((\mathrm C_2)\)}
Condition \((\mathrm C_0)\) holds, and for every \(S'\) over \(S\),
every quasi-regular section of
\[
  \mathfrak g_{S'}=\mathfrak g\otimes_S S'
\]
is regular.

\smallskip
\noindent
\textup{\((\mathrm C_3)\)}
Every \(s\in S\) has an open neighborhood \(V\) on which the
Killing polynomial of \(\mathfrak g\) has the form
\[
  t^r\bigl(t^{n-r}+c_1t^{n-r-1}+\cdots+c_{n-r}\bigr),
\]
where, for every \(s\in V\),
\[
  c_{n-r}(s)\in
  \operatorname{Sym}\!\bigl(\mathfrak g(s)^\vee\bigr)
\]
is nonzero.

\smallskip
With this notation, one has the implications
\[
  (\mathrm C_3)\Longrightarrow(\mathrm C_2)
  \Longrightarrow(\mathrm C'_1)\Longrightarrow(\mathrm C_1)
  \Longrightarrow(\mathrm C_0),
\]
and, when \(S\) is reduced, the five conditions above are
equivalent.}

Notice also that these conditions are manifestly stable under base
change and are local for the fpqc topology.

The implications
\[
  (\mathrm C'_1)\Longrightarrow(\mathrm C_1)
  \Longrightarrow(\mathrm C_0)
\]
are trivial, and the implication
\((\mathrm C_0)\Longrightarrow(\mathrm C_3)\) when \(S\) is
reduced is immediate. We note moreover:

\paragraph{\textbf{Corollary 2.10.}}
\SGARefTarget{sga3:xiv:2.10}
\emph{Suppose that condition \((\mathrm C_0)\) holds. Let \(U\) be
the set of elements of \(W(\mathfrak g)\) that are regular in their
fiber. Then \(U\) is open. In particular, for every section \(a\)
of \(\mathfrak g\) over \(S\), the set \(V\) of \(s\in S\) for
which \(a(s)\in\mathfrak g(s)\) is regular is open.}

Indeed, the first assertion follows from the second, applied to
\(\mathfrak g_{S'}\) for every base change \(S'\) over \(S\). For
the second, we may suppose that \(S\) is reduced; then
\((\mathrm C_3)\) holds, and it is enough to inspect the Killing
polynomial of \(a\) in \(\mathfrak g\).

The implication
\((\mathrm C_2)\Longrightarrow(\mathrm C'_1)\) now follows easily
from part~b) of the following more precise statement.

\paragraph{\textbf{Corollary 2.11.}}
\SGARefTarget{sga3:xiv:2.11}
{\itshape Suppose that condition \((\mathrm C_2)\) holds. Then:

\smallskip
\noindent
\textup{a)}
For every \(s\in S\) and every Cartan subalgebra
\(\mathfrak d_0\) of \(\mathfrak g(s)\) such that
\(\mathfrak d_0\) contains a regular element of \(\mathfrak g(s)\)
\textup{(a condition automatically satisfied when \(k(s)\) is
infinite)}, there exist an open neighborhood \(V\) of \(s\) and a
Cartan subalgebra \(\mathfrak d\) of \(\mathfrak g|_V\) whose fiber
at \(s\) is \(\mathfrak d_0\). If \(S_1\) is a subscheme of \(S\)
containing \(s\), and if \(\mathfrak d_0\) has already been
extended to a Cartan subalgebra \(\mathfrak d_1\) of
\(\mathfrak g\otimes_S S_1\), then one can find an open
neighborhood \(V\) of \(s\) in \(S\) and a Cartan subalgebra
\(\mathfrak d\) of \(\mathfrak g|_V\) such that
\[
  \mathfrak d\otimes_V(S_1\cap V)
  =\mathfrak d_1|_{S_1\cap V}.
\]

\smallskip
\noindent
\textup{b)}
For every \(s\in S\) such that \(\mathfrak g(s)\) contains a
regular element \textup{(a condition automatically satisfied when
\(k(s)\) is infinite)}, there exist an open neighborhood \(V\) of
\(s\) and a Cartan subalgebra \(\mathfrak d\) of
\(\mathfrak g|_V\).}

Part~b) follows from part~a) by taking
\[
  \mathfrak d_0=\Nil\!\left(a_0,\mathfrak g(s)\right),
\]
where \(a_0\) is a regular element of \(\mathfrak g(s)\).
To prove part~a), in its second formulation, choose a regular
element \(a_0\) of \(\mathfrak g(s)\) contained in
\(\mathfrak d_0\). Extend it near \(s\) in \(S_1\) to a section of
\(\mathfrak d_1\), and extend that section in turn near \(s\) to a
section of \(\mathfrak g\). By \SGARefLink{sga3:xiv:2.10}{2.10} this section is quasi-regular
on an open neighborhood \(V\) of \(s\), hence regular by
\((\mathrm C_2)\). Therefore, by \SGARefLink{sga3:xiv:2.6}{2.6},
\[
  \mathfrak d=\Nil\!\left(a,\mathfrak g|_V\right)
\]
has the required properties.

% END inputs/XIV/tex/components/04_expose_XIV_proposition29_through_corollary211_en.tex

% BEGIN inputs/XIV/tex/components/05_expose_XIV_remark212_lemma213_en.tex
% Authority: Expo14.pdf, local p. 7 / combined p. 841.
% Continuous English ends after the proof of Lemma 2.13, hence of 2.9.

\paragraph{\textbf{Remark 2.12.}}
\SGARefTarget{sga3:xiv:2.12}
I do not know whether the assertions of \SGARefLink{sga3:xiv:2.11}{2.11}~a) and b) remain valid
without the hypothesis that regular points exist, when \(k(s)\) is
finite.

It remains to prove
\[
  (\mathrm C_3)\Longrightarrow(\mathrm C_2).
\]
We also record the following equivalent form of
\((\mathrm C_3)\):

\smallskip
\noindent
\textup{\((\mathrm C'_3)\)}
Condition \((\mathrm C_0)\) holds, that is, the nilpotent rank of
\(\mathfrak g(s)\), for \(s\in S\), is locally constant; and on
every open subset \(V\) of \(S\) on which this rank has the value
\(r\), the Killing polynomial of \(\mathfrak g|_V\) is divisible
by \(t^r\).

\smallskip
We must show that this condition implies that every quasi-regular
section \(a\) of \(\mathfrak g\) is regular. In view of \SGARefLink{sga3:xiv:2.6}{2.6}, this
follows from the implication
\(\textup{(iv)}\Longrightarrow\textup{(iii)}\) in the following
lemma, applied to the endomorphism
\(\operatorname{ad}(a)\) of \(\mathfrak g\).

\paragraph{\textbf{Lemma 2.13.}}
\SGARefTarget{sga3:xiv:2.13}
{\itshape Let \(A\) be a ring, let \(M\) be a projective
\(A\)-module of finite type, and let \(u\) be an endomorphism of
\(M\). The following conditions are equivalent:

\smallskip
\noindent
\textup{(i)}
\(M\) is the direct sum of two \(u\)-stable submodules \(M'\) and
\(M''\), such that \(u|_{M'}\) is nilpotent and \(u|_{M''}\) is an
automorphism of \(M''\).

\smallskip
\noindent
\textup{(ii)}
There exists an integer \(n>0\) such that
\[
  \operatorname{Im}u^n+\operatorname{Ker}u^n=M.
\]

\smallskip
\noindent
\textup{(iii)}
The nilspace
\[
  \Nil(u)=\bigcup_{n>0}\operatorname{Ker}u^n
\]
is a direct factor of \(M\), and
\[
  M=\Nil(u)+u(M).
\]

\smallskip
These conditions are implied by the following one, and are
equivalent to it when \(A\) is reduced:

\smallskip
\noindent
\textup{(iv)}
Locally on \(\operatorname{Spec}(A)\), for the Zariski topology, the
characteristic polynomial \(P_u(t)\) of \(u\) can be written in the
form
\[
  t^r\bigl(t^{n-r}+c_1t^{n-r-1}+\cdots+c_{n-r}\bigr),
\]
where \(c_{n-r}\) is invertible.}

The equivalence of \textup{(i)}, \textup{(ii)}, and \textup{(iii)}
is immediate and is recalled for reference. When \(A\) is reduced,
\textup{(i)} implies \textup{(iv)} because a nilpotent
endomorphism of a projective module of rank \(r\) then has
characteristic polynomial \(t^r\). On the other hand, in all cases
the constant term of the characteristic polynomial of an
automorphism of a projective module of finite type is, up to sign,
the determinant of \(u\), locally on \(\operatorname{Spec}(A)\);
it is therefore an invertible element of \(A\).

Finally, to prove
\(\textup{(iv)}\Longrightarrow\textup{(i)}\), regard \(M\) as a
module over the polynomial ring \(A[t]\), with \(t\) acting as
\(u\). The familiar identity
\[
  P(u)=0
\]
shows that \(M\) is annihilated by \(P A[t]\), and hence may be
regarded as a module over \(A[t]/P A[t]\). Write
\[
  P=t^rQ,
\]
where the constant term of \(Q\) is invertible. Then
\[
  P A[t]=t^rA[t]\cap Q A[t],
\]
so that \(A[t]/P A[t]\) decomposes as the product of the rings
\[
  A[t]/t^rA[t]
  \qquad\text{and}\qquad
  A[t]/Q A[t].
\]
This gives a corresponding decomposition of \(M\) into the direct
sum of two \(A[t]\)-modules, or equivalently two \(u\)-stable
\(A\)-submodules \(M'\) and \(M''\); this is precisely the
decomposition of \textup{(i)}.

This completes the proof of \SGARefLink{sga3:xiv:2.13}{2.13}, and hence that of \SGARefLink{sga3:xiv:2.9}{2.9}.

% END inputs/XIV/tex/components/05_expose_XIV_remark212_lemma213_en.tex

% BEGIN inputs/XIV/tex/components/06_expose_XIV_corollary214_remark215_en.tex
% Authority: Expo14.pdf, local p. 8 / combined p. 842.
% Continuous English ends immediately after Remark 2.15(c).

\paragraph{\textbf{Corollary 2.14.}}
\SGARefTarget{sga3:xiv:2.14}
\emph{When condition \((\mathrm C_3)\) of \SGARefLink{sga3:xiv:2.9}{2.9} holds, the Cartan
subalgebras of \(\mathfrak g\) are strictly nilpotent
\textup{(\SGARefLink{sga3:xiv:2.2}{2.2})}.}

The proof is immediate.

\paragraph{\textbf{Remark 2.15.}}
\SGARefTarget{sga3:xiv:2.15}
\noindent
\textup{a)}
One can prove a converse to \SGARefLink{sga3:xiv:2.14}{2.14}: condition \((\mathrm C_3)\) is
equivalent to the assertion that, for every \(S'\) over \(S\),
every quasi-regular section of \(\mathfrak g_{S'}\) is regular and
every Cartan subalgebra of \(\mathfrak g_{S'}\) is strictly
nilpotent; equivalently, every quasi-regular section of
\(\mathfrak g_{S'}\) is contained in a strictly nilpotent Cartan
subalgebra of \(\mathfrak g_{S'}\).

\smallskip
\noindent
\textup{b)}
Unlike the other conditions \((\mathrm C_0)\) through
\((\mathrm C_2)\), condition \((\mathrm C_3)\) is of
``infinitesimal nature.'' More precisely, when \(S\) is locally
noetherian, \(\mathfrak g\) satisfies \((\mathrm C_3)\) if and only
if, for every base change \(S'\to S\) with \(S'\) local Artinian,
\(\mathfrak g_{S'}\) satisfies the same condition. If desired, one
may take \(S'\) to be the spectrum of an Artinian quotient of a
local ring of \(S\). Similarly, when \((\mathrm C_0)\) holds, the
condition for a section of \(\mathfrak g\) to be regular is of
infinitesimal nature.

\smallskip
\noindent
\textup{c)}
When \(\mathfrak g\) is the Lie algebra of a smooth group
prescheme of finite presentation over \(S\), we shall see that the
conditions
\[
  (\mathrm C_0),\quad(\mathrm C_1),\quad
  (\mathrm C'_1),\quad(\mathrm C_2)
\]
on \(\mathfrak g\) are equivalent \textup{(\SGARefLink{sga3:xiv:5.2}{5.2}~a)}. I do not know
what happens in general, except that even when \(S\) is local
Artinian, \((\mathrm C_0)\) does not imply \((\mathrm C_1)\).
However, even when \(\mathfrak g\) comes from a group \(G\) and
\(S\) is local Artinian, it is not true in general that
\((\mathrm C_2)\) implies \((\mathrm C_3)\), since
\(\mathfrak g\) may be nilpotent without being strictly
nilpotent.

An example is obtained from a smooth affine group scheme \(G\)
over the spectrum \(S\) of a discrete valuation ring, such that
the Lie algebra of the generic fiber is not nilpotent---for
example, the generic fiber may be an adjoint semisimple
group---whereas that of the special fiber is nilpotent---for
example, the special fiber may be a vector group. Then, for
sufficiently large \(n\), the Lie algebra of
\[
  G_n=G\times_S S_n
\]
is not strictly nilpotent, although it is nilpotent.

% END inputs/XIV/tex/components/06_expose_XIV_corollary214_remark215_en.tex

% BEGIN inputs/XIV/tex/components/07_expose_XIV_cartans_functor_theorem216_statement_en.tex
% Authority: Expo14.pdf, local pp. 8--9 / combined pp. 842--843.
% Continuous English ends immediately before the proof of Theorem 2.16.

We continue under the hypotheses of \SGARefLink{sga3:xiv:2.4}{2.4}. Let
\[
  \mathcal D\colon
  \bigl((\mathrm{Sch})_{/S}\bigr)^\circ
  \longrightarrow \mathrm{Sets}
\]
be the functor defined by
\[
  \mathcal D(S')
  =\left\{
    \text{Cartan subalgebras of }\mathfrak g_{S'}
   \right\}.
\]
Introduce also the functor
\[
  X(S')
  =\left\{
    (\mathfrak d,a)\ \middle|\
    \begin{array}{l}
      \mathfrak d\text{ is a Cartan subalgebra of }\mathfrak g_{S'},\\
      a\text{ is a section of }\mathfrak d
    \end{array}
   \right\}.
\]
The two projections
\[
  (\mathfrak d,a)\longmapsto\mathfrak d
  \qquad\text{and}\qquad
  (\mathfrak d,a)\longmapsto a
\]
give morphisms
\[
  p\colon X\longrightarrow\mathcal D
  \qquad\text{and}\qquad
  \psi\colon X\longrightarrow W(\mathfrak g).
\]

When \((\mathrm C_0)\) holds, let \(U\) also denote the open subset
of regular points of \(W(\mathfrak g)\) \textup{(\SGARefLink{sga3:xiv:2.10}{2.10})}. Condition
\((\mathrm C_2)\) is then equivalent to saying that the morphism
\[
  \psi^{-1}(U)\longrightarrow U
\]
induced by \(\psi\) is an isomorphism; a priori it is a
monomorphism by \SGARefLink{sga3:xiv:2.6}{2.6}.

It is immediate that
\[
  p\colon X\longrightarrow\mathcal D
\]
is representable by a projection of vector bundles. More
explicitly, for every \(S\)-morphism \(S'\to\mathcal D\)
corresponding to a Cartan subalgebra \(\mathfrak d\) of
\(\mathfrak g_{S'}\), the fiber product
\[
  X\times_{\mathcal D}S'
\]
is represented by the vector bundle \(W(\mathfrak d)\) over
\(S'\). Thus, if \(\mathcal D\) is representable, then so is \(X\).
We have:

\paragraph{\textbf{Theorem 2.16.}}
\SGARefTarget{sga3:xiv:2.16}
{\itshape Let \(S\) be a prescheme and let \(\mathfrak g\) be a Lie
algebra over \(S\) whose underlying \(\mathcal O_S\)-module is
locally free of finite type. Suppose that condition
\((\mathrm C_0)\) of \SGARefLink{sga3:xiv:2.9}{2.9} holds.

\smallskip
\noindent
\textup{a)}
The functor \(\mathcal D\) of Cartan subalgebras of
\(\mathfrak g\), defined above, is represented by a
quasi-projective prescheme of finite presentation over \(S\).
The same holds for the functor \(X\).

\smallskip
\noindent
\textup{b)}
If condition \((\mathrm C_2)\) of \SGARefLink{sga3:xiv:2.9}{2.9} holds, then
\(\mathcal D\) and \(X\) are smooth over \(S\), and the morphism
\[
  \psi^{-1}(U)\longrightarrow U
\]
induced by \(\psi\) is an isomorphism.

\smallskip
\noindent
\textup{c)}
Continue to suppose that \((\mathrm C_2)\) holds. Let \(s\in S\),
and let \(\mathfrak d_0\) be a Cartan subalgebra of
\(\mathfrak g(s)\), corresponding to a \(k(s)\)-rational point
\(d\) of \(\mathcal D\). Suppose that \(\mathfrak d_0\) contains a
regular element of \(\mathfrak g(s)\), a condition automatically
satisfied when \(k(s)\) is infinite. If \(r\) is the infinitesimal
rank of \(\mathfrak g(s)\) and \(n\) its rank over \(k(s)\), then
there exists an open neighborhood \(V\) of \(d\) in
\(\mathcal D\) that is \(S\)-isomorphic to an open subset \(V'\)
of
\[
  S[t_1,\ldots,t_{n-r}].
\]}

% END inputs/XIV/tex/components/07_expose_XIV_cartans_functor_theorem216_statement_en.tex

% BEGIN inputs/XIV/tex/components/08_expose_XIV_theorem216_proof_corollary217_remark218_en.tex
% Authority: Expo14.pdf, local pp. 9--11 / combined pp. 843--845.
% Continuous English ends at the close of Section 2, before Section 3.

\paragraph{\textit{Proof.}}
We may suppose that \(\mathfrak g\) has constant rank \(n\) and
constant nilpotent rank \(r\). The assertions concerning \(X\) in
\textup{a)} and \textup{b)} follow immediately from those concerning
\(\mathcal D\) and from the fact that \(X\) is a vector bundle over
\(\mathcal D\), defined by a locally free module; they were included
only for completeness.

\medskip
\noindent\textup{a)}
The functor \(\mathcal D\) is a subfunctor of
\(\operatorname{Grass}_{n-r}(\mathfrak g)\), whose value at \(S'\)
is the set of locally free quotient modules of rank \(n-r\) of
\(\mathfrak g_{S'}\). It is well known that the latter functor is
represented by a projective prescheme smooth over \(S\); see, for
example, Cartan Seminar 1960/61, Exposé~12, Nos.~2 and~3, whose
constructions carry over unchanged to preschemes.\begingroup
\renewcommand{\thefootnote}{*}\footnote{See also EGA~I, second
edition \textup{(announced as forthcoming from North-Holland
Publishing Company)}.}
\endgroup\footnote{Editors' note: See \S \SGARefLink{sga3:xiv:1.3}{I.1.3} of M.~Demazure and
P.~Gabriel, \emph{Groupes algébriques}, Masson \textup{(1970)}.}

We are therefore reduced to the following relative problem. Given
a locally free quotient module of rank \(n-r\) of \(\mathfrak g\),
or, equivalently, a locally free submodule \(\mathfrak d\) of rank
\(r\) that is locally a direct factor, represent the functor
\[
 F(S')=
 \begin{cases}
   \varnothing,
     &\text{if \(\mathfrak d_{S'}\) is not a Cartan subalgebra of
       \(\mathfrak g_{S'}\),}\\
   \{\varnothing\},
     &\text{otherwise.}
 \end{cases}
\]
We shall show that \(F\) is represented by a subprescheme of finite
presentation over \(S\). This will show that
\[
  \mathcal D\longrightarrow
  \operatorname{Grass}_{n-r}(\mathfrak g)
\]
is represented by an immersion of finite presentation and will
complete the proof of \textup{a)}.

First express the condition that \(\mathfrak d_{S'}\) be a Lie
subalgebra of \(\mathfrak g_{S'}\). This condition says precisely
that \(S'\to S\) factors through a certain closed subprescheme
\(S_1\) of finite presentation over \(S\); locally on \(S\), its
equations are obtained at once from a basis of \(\mathfrak g\)
adapted to \(\mathfrak d\). We may therefore suppose that
\(S=S_1\).

Next we must express that \(\mathfrak d_{S'}\) contains,
fpqc-locally, a quasi-regular section of \(\mathfrak g_{S'}\). For
this purpose consider
\[
  V=W(\mathfrak d)\cap U,
\]
where \(U\) is the open subset of regular points of
\(W(\mathfrak g)\) \textup{(\SGARefLink{sga3:xiv:2.10}{2.10})}. Since the structure morphism
\(V\to S\) is smooth and quasi-compact, its image \(S_1\) is an
open subset of \(S\), and the immersion \(S_1\to S\) is
quasi-compact, hence of finite presentation. The condition on
\(S'\) is exactly that \(S'\to S\) factor through \(S_1\). We are
thus reduced to the case \(S=S_1\), and, by descent, to the case in
which \(\mathfrak d\) has a section \(a\) that is quasi-regular in
\(\mathfrak g\).

It remains to express that, for the induced section \(a_{S'}\), the
endomorphism
\[
  \operatorname{ad}(a_{S'})
    _{\mathfrak g_{S'}/\mathfrak d_{S'}}
\]
is bijective. This again says that \(S'\to S\) factors through an
open subprescheme of finite presentation, namely \(S_D\), where
\[
  D=\det\!\left(
    \operatorname{ad}(a)_{\mathfrak g/\mathfrak d}
  \right).
\]
But \(\mathfrak d|_{S_D}\) is then a Cartan subalgebra of
\(\mathfrak g|_{S_D}\). Consequently \(S_D\) represents \(F\),
which proves \textup{a)}.

\medskip
\noindent\textup{b)}
This follows immediately from \SGARefLink{sga3:xiv:2.11}{2.11}\textup{(a)} and \SGARefLink{sga3:xi:remark:1:5}{XI~1.5}. It is
also, of course, a consequence of the more precise assertion
\textup{c)}.

\medskip
\noindent\textup{c)}
Let \(a_0\) be a regular point of \(\mathfrak g(s)\) contained in
\(\mathfrak d_0\), and extend it to a section \(a\) of
\(\mathfrak g\) over an open neighborhood of \(s\). We may plainly
suppose that this neighborhood is all of \(S\). Choose a complement
\(M_0\) of \(\mathfrak d_0\) in the vector space
\(\mathfrak g(s)\). After again replacing \(S\) by an open
neighborhood of \(s\), there is a submodule \(M\) of
\(\mathfrak g\), locally a direct factor, such that
\(M(s)=M_0\); we again suppose that the neighborhood is \(S\).

Let \(\mathcal V\) be the subfunctor of \(\mathcal D\) for which
\(\mathcal V(S')\) consists of the Cartan subalgebras
\(\mathfrak d'\) of \(\mathfrak g_{S'}\) satisfying:
\begin{enumerate}
\item[\textup{1\textdegree)}]
  \(\mathfrak d'\) is a complement of \(M_{S'}\);
\item[\textup{2\textdegree)}]
  the unique section in
  \((a_{S'}+M_{S'})\cap\mathfrak d'\) is a regular section of
  \(\mathfrak g_{S'}\).
\end{enumerate}
Condition \textup{1\textdegree)} defines an open \(V_1\) of
\(\mathcal D\), induced by the open of
\(\operatorname{Grass}_{n-r}(\mathfrak g)\) given by the same
condition. By \SGARefLink{sga3:xiv:2.10}{2.10} and \((\mathrm C_2)\), the conjunction of
\textup{1\textdegree)} and \textup{2\textdegree)} defines an open
of \(V_1\). Thus \(\mathcal V\) is represented by an open
subprescheme \(V\) of \(\mathcal D\), plainly containing the point
\(d\).

Now let \(\mathcal V'\) be the subfunctor of \(W(M)\) defined by
\[
 \mathcal V'(S')
 =
 \left\{
   u'\in\Gamma(S',M_{S'})
   \ \middle|\
   \begin{array}{l}
   \textup{(i) }a_{S'}+u'\text{ is a regular section of }
     \mathfrak g_{S'},\\
   \textup{(ii) the unique Cartan subalgebra }\mathfrak d'
     \text{ of }\mathfrak g_{S'}\\
   \quad\text{containing }a_{S'}+u'
     \text{ is a complement of }M_{S'}
   \end{array}
 \right\}.
\]
Condition \textup{(i)} defines an open subprescheme \(V'_1\) of
\(W(M)\): it is the inverse image of the open \(U\) of regular
points of \(W(\mathfrak g)\) \textup{(\SGARefLink{sga3:xiv:2.10}{2.10})} under the translation
morphism \(m\mapsto a+m\). The conjunction of \textup{(i)} and
\textup{(ii)} defines an open \(V'\) of \(V'_1\), namely the
inverse image of \(V\) under the evident morphism
\[
  V'_1\longrightarrow\mathcal D
\]
that associates with \(u'\) the unique Cartan subalgebra
\(\mathfrak d'\) containing \(a_{S'}+u'\). Its restriction gives
a morphism
\[
  V'\longrightarrow V,
\]
which is plainly an isomorphism. This proves \textup{c)}.

\paragraph{\textbf{Corollary 2.17.}}
\SGARefTarget{sga3:xiv:2.17}
{\itshape Let \(\mathfrak g\) be a finite-dimensional Lie algebra
over a field \(k\). Then the scheme \(\mathcal D\) of Cartan
subalgebras of \(\mathfrak g\) \textup{(\SGARefLink{sga3:xiv:2.16}{2.16}\textup{(a)})} is
quasi-projective, smooth, and irreducible. If \(\mathfrak g\)
contains a regular element \textup{(for example, if \(k\) is
infinite)}, then \(\mathcal D\) is a rational variety; that is, its
function field is a purely transcendental extension of \(k\).}

Indeed, \(\mathcal D\) is irreducible because there is a surjective
morphism
\[
  \psi^{-1}(U)\longrightarrow\mathcal D,
\]
and \(\psi^{-1}(U)\) is irreducible, being isomorphic to the open
subset \(U\) of \(W(\mathfrak g)\). The assertion concerning the
function field follows immediately from \SGARefLink{sga3:xiv:2.16}{2.16}\textup{(c)}.

\paragraph{\textbf{Remarks 2.18.}}
\SGARefTarget{sga3:xiv:2.18}
I do not know whether this conclusion remains true when \(k\) is
finite without assuming that \(\mathfrak g\) contains a regular
point; compare \SGARefLink{sga3:xiv:2.12}{2.12}. It can be proved when \(\mathfrak g\) is the
Lie algebra of an algebraic group \(G\) smooth over \(k\), at least
when \(G/\operatorname{radical}(G)\) is an ``adjoint'' semisimple
group, by using a result of Chevalley cited below \textup{(see the
Appendix)}. It is plausible that the result remains valid without
restriction on \(G\); for this it would suffice that the cited
result of Chevalley be proved for every semisimple algebraic group,
not necessarily an adjoint group.

% END inputs/XIV/tex/components/08_expose_XIV_theorem216_proof_corollary217_remark218_en.tex

% BEGIN inputs/XIV/tex/components/09_expose_XIV_theorem31_en.tex
% Authority: Expo14.pdf, local pp. 11--12 / combined pp. 845--846.
% Continuous English ends after the proof of Theorem 3.1.

\section{Subgroups of type \textup{(C)} of group preschemes over an
arbitrary prescheme}
\SGARefTarget{sga3:xiv:section:3}

\paragraph{\textbf{Theorem 3.1.}}
\SGARefTarget{sga3:xiv:3.1}
{\itshape Let \(S\) be a prescheme, let \(G\) be a smooth
\(S\)-group prescheme, and let \(\mathfrak g\) be its Lie algebra,
which is a locally free module of finite type over \(S\). Let
\(\mathfrak h\) be a Lie subalgebra of \(\mathfrak g\) that, as a
module, is locally a direct factor in \(\mathfrak g\), and suppose
that, for every \(s\in S\), the geometric fiber
\(\mathfrak h_{\bar s}\) contains a Cartan subalgebra of
\(\mathfrak g_{\bar s}\). Let \(a\) be a quasi-regular section of
\(\mathfrak g\) \textup{(\SGARefLink{sga3:xiv:2.5}{2.5})}. Then
\[
  M_a=\operatorname{Transp}_G(a,\mathfrak h),
\]
the subfunctor of \(G\) whose \(S'\)-valued points are the
\(g\in G(S')\) such that
\[
  \operatorname{ad}(g)\cdot a_{S'}
    \in\Gamma(S',\mathfrak h_{S'}),
\]
is represented by a closed subprescheme of \(G\), smooth over
\(S\), whose structure morphism to \(S\) is surjective.}
\footnote{Editors' note: This transporter is denoted
\(\operatorname{Transp}_G(a,\mathfrak h)\).}

\paragraph{\textit{Proof.}}
Consider the canonical morphism
\[
  \varphi\colon G\times_S W(\mathfrak h)\longrightarrow
  W(\mathfrak g),
  \qquad
  (g,x)\longmapsto\operatorname{ad}(g)\cdot x.
\]
Then \(M_a\) is \(S\)-isomorphic to \(\varphi^{-1}(a)\), the inverse
image under \(\varphi\) of \(a\), viewed as a section of
\(W(\mathfrak g)\) over \(S\). To prove that \(M_a\) is smooth, it
is enough to show that \(\varphi\) is smooth at the points of
\(G\times_S W(\mathfrak h)\) lying over \(\operatorname{Im}(a)\).
More generally, \(\varphi\) is smooth at every point lying over a
point of \(W(\mathfrak g)\) that is regular in its fiber
\(W(\mathfrak g(s))\) over \(S\).

Indeed, the source and target of \(\varphi\) are smooth over \(S\),
and hence are flat and locally of finite presentation over \(S\).
We may therefore check the assertion fiber by fiber. This reduces
us to the case in which \(S\) is the spectrum of an algebraically
closed field \(k\), \(G\) is a locally algebraic group over \(k\),
\(\mathfrak h\) is a Lie subalgebra of its Lie algebra
\(\mathfrak g\) containing a Cartan subalgebra of \(\mathfrak g\),
and \(a\) is a regular point of \(\mathfrak g\).

Since \(\varphi\) is a \(G\)-morphism, we may suppose that the point
of \(G\times W(\mathfrak h)\) under consideration is of the form
\((e,a)\). We may also suppose that \(G\) is connected, and hence
of finite type over \(k\). The assertion is then precisely
\SGARefLink{sga3:xiii:corollary:5:4}{XIII~5.4}.

It is also immediate that \(M_a\) is a closed subprescheme of \(G\)
of finite presentation over \(S\): it is the inverse image of
\(W(\mathfrak h)\) under the morphism
\[
  G\longrightarrow W(\mathfrak g),
  \qquad
  g\longmapsto\operatorname{ad}(g)\cdot a.
\]
Finally, the surjectivity of \(M_a\to S\) likewise reduces to the
case of an algebraically closed base field. By hypothesis,
\(\mathfrak h\) then contains a Cartan subalgebra
\(\mathfrak d'\). By the conjugacy theorem \SGARefLink{sga3:xiii:theorem:6:1}{XIII~6.1}\textup{(a)},
\(\mathfrak d'\) is conjugate to
\[
  \mathfrak d=\Nil(a,\mathfrak g).
\]
It follows that \(\mathfrak h\) contains a conjugate of \(a\).
This completes the proof.

% END inputs/XIV/tex/components/09_expose_XIV_theorem31_en.tex

% BEGIN inputs/XIV/tex/components/10_expose_XIV_corollary32_lemma33_en.tex
% Authority: Expo14.pdf, local p. 12 / combined p. 846.
% Continuous English ends immediately before Corollary 3.4.

\paragraph{\textbf{Corollary 3.2.}}
\SGARefTarget{sga3:xiv:3.2}
{\itshape Let \(G\) and \(\mathfrak g\) be as in \SGARefLink{sga3:xiv:3.1}{3.1}, with \(G\)
of finite type over \(S\), and let \(\mathfrak k\) and
\(\mathfrak h\) be two Lie subalgebras of \(\mathfrak g\) that are
locally direct factors as modules. Suppose that one of the following
two hypotheses holds:

\smallskip
\noindent
\textup{a)}
For every \(s\in S\), the geometric fiber
\(\mathfrak k_{\bar s}\) is nilpotent and contains a regular element
of \(\mathfrak g_{\bar s}\), while the geometric fiber
\(\mathfrak h_{\bar s}\) contains a Cartan subalgebra of
\(\mathfrak g_{\bar s}\).

\smallskip
\noindent
\textup{b)}
\(\mathfrak k\) is a Cartan subalgebra of \(\mathfrak g\).

\smallskip
\noindent
Then \(\operatorname{Transp}_G(\mathfrak k,\mathfrak h)\) is a
closed subprescheme of \(G\), smooth over \(S\). Moreover, in case
\textup{a)}, its structure morphism to \(S\) is surjective.}

The representability of the transporter by a closed subprescheme of
\(G\) of finite presentation is immediate and is left to the
reader. We first prove smoothness in case \textup{a)}. Suppose first
that \(\mathfrak k\) has a section \(a\) that is quasi-regular in
\(\mathfrak g\). It is then enough to apply \SGARefLink{sga3:xiv:3.1}{3.1} and the following
lemma.

\paragraph{\textbf{Lemma 3.3.}}
\SGARefTarget{sga3:xiv:3.3}
{\itshape Under the hypotheses of \SGARefLink{sga3:xiv:3.2}{3.2}\textup{(a)}, if \(a\) is a
section of \(\mathfrak k\) that is quasi-regular in
\(\mathfrak g\), then
\[
  \operatorname{Transp}_G(\mathfrak k,\mathfrak h)
  =
  \operatorname{Transp}_G(a,\mathfrak h).
\]}

Indeed, by the definitions it is enough to show that, if \(a\) is
also a section of \(\mathfrak h\), then
\(\mathfrak k\subset\mathfrak h\). Since \(\mathfrak k\) is
locally nilpotent by hypothesis, \SGARefLink{sga3:xiv:2.1}{2.1} gives
\[
  \mathfrak k\subset\Nil(a,\mathfrak g).
\]
On the other hand,
\[
  \Nil(a,\mathfrak g)\subset\mathfrak h,
\]
because
\[
  \operatorname{ad}(a)_{\mathfrak g/\mathfrak h}
\]
is injective; this can be checked fiber by fiber by \SGARefLink{sga3:xiii:proposition:4:8:item:b}{XIII~4.8(b)}.
The lemma follows.

In the general case, the usual standard argument first reduces us
to the case where \(S\) is affine noetherian, and then to the case
where \(S\) is local artinian, smoothness being an infinitesimal
property. By flat descent we may further suppose that the residue
field is infinite. The fiber of \(\mathfrak k\) then contains an
element regular in the corresponding fiber of \(\mathfrak g\).
Lifting this element to a section of \(\mathfrak k\) reduces us to
the preceding case.

Thus the transporter is smooth in case \textup{a)}. Its structure
morphism is surjective because this may be checked after reducing to
the case where \(S\) is the spectrum of an algebraically closed
field; \(\mathfrak k\) then contains a regular point of
\(\mathfrak g\), and \SGARefLink{sga3:xiv:3.3}{3.3} and \SGARefLink{sga3:xiv:3.1}{3.1} apply.

For case \textup{b)}, the definition of smoothness \textup{(\SGARefLink{sga3:xi:definition:1:1}{XI~1.1})}
reduces the assertion to the following lifting statement. Let \(S\)
be affine, let \(S_0\) be a closed subscheme defined by a
quasi-coherent nilpotent ideal \(J\), and let
\(g_0\in G(S_0)\) transport \(\mathfrak k_0\) into
\(\mathfrak h_0\). Then \(g_0\) lifts to an element \(g\in G(S)\)
that transports \(\mathfrak k\) into \(\mathfrak h\). The
hypothesis on \(g_0\) places this lifting problem under case
\textup{a)}, already proved. This completes the proof.

Of course, if in \SGARefLink{sga3:xiv:3.2}{3.2}\textup{(b)} the subalgebra \(\mathfrak h\)
satisfies the stronger hypothesis of \SGARefLink{sga3:xiv:3.1}{3.1}, then, and only then, the
structure morphism
\[
  \operatorname{Transp}_G(\mathfrak k,\mathfrak h)
  \longrightarrow S
\]
is surjective. Using Hensel's lemma \SGARefLink{sga3:xi:proposition:1:10}{XI~1.10}, one obtains from \SGARefLink{sga3:xiv:3.1}{3.1}
and \SGARefLink{sga3:xiv:3.2}{3.2} the following corollary.

% END inputs/XIV/tex/components/10_expose_XIV_corollary32_lemma33_en.tex

% BEGIN inputs/XIV/tex/components/11_expose_XIV_corollary34_through_definition38_en.tex
% Authority: Expo14.pdf, local pp. 12--13 / combined pp. 846--847.
% Continuous English ends immediately before Theorem 3.9.

\paragraph{\textbf{Corollary 3.4.}}
\SGARefTarget{sga3:xiv:3.4}
{\itshape Under the hypotheses of \SGARefLink{sga3:xiv:3.1}{3.1} for \(G\) and
\(\mathfrak h\), and assuming that \(G\) is of finite type over
\(S\), the following assertions hold:

\smallskip
\noindent
\textup{a)}
For every quasi-regular section \(a\) of \(\mathfrak g\), étale
locally there is a conjugate of \(a\) that is a section of
\(\mathfrak h\).

\smallskip
\noindent
\textup{b)}
For every Cartan subalgebra \(\mathfrak d\) of \(\mathfrak g\),
étale locally \(\mathfrak d\) is conjugate to a subalgebra of
\(\mathfrak h\).}

In particular, when \(\mathfrak h\) is itself a Cartan subalgebra
of \(\mathfrak g\), one obtains the following.

\paragraph{\textbf{Corollary 3.5.}}
\SGARefTarget{sga3:xiv:3.5}
{\itshape Let \(G\) be a smooth \(S\)-prescheme in groups of finite
type, let \(\mathfrak g\) be its Lie algebra, and let
\(\mathfrak d\) and \(\mathfrak d'\) be two Cartan subalgebras of
\(\mathfrak g\). Then
\(\operatorname{Transp}_G(\mathfrak d,\mathfrak d')\) is identical
to the strict transporter of \(\mathfrak d\) into
\(\mathfrak d'\), and is a closed subprescheme of \(G\), smooth
over \(S\), whose structure morphism is surjective. Étale locally,
\(\mathfrak d\) and \(\mathfrak d'\) are conjugate.}

Here the transporter is identical to the strict transporter simply
because \(\mathfrak d\) and \(\mathfrak d'\) are locally direct
factors in \(\mathfrak g\) and have the same rank at every point.
Thus \SGARefLink{sga3:xiv:3.5}{3.5} is a special case of \SGARefLink{sga3:xiv:3.4}{3.4}. Taking
\(\mathfrak d=\mathfrak d'\) gives the following.

\paragraph{\textbf{Corollary 3.6.}}
\SGARefTarget{sga3:xiv:3.6}
{\itshape Let \(G\) be as in \SGARefLink{sga3:xiv:3.5}{3.5}, and let \(\mathfrak d\) be a
Cartan subalgebra of \(\mathfrak g\). Then
\(\operatorname{Norm}_G(\mathfrak d)\) is a closed subprescheme in
groups of \(G\), smooth over \(S\), and its Lie algebra is
\(\mathfrak d\).}

Indeed, the last assertion says that \(\mathfrak d\) is its own
normalizer in \(\mathfrak g\), and this follows at once because it
is true fiber by fiber.

\paragraph{\textbf{Corollary 3.7.}}
\SGARefTarget{sga3:xiv:3.7}
{\itshape Let \(G\) and \(\mathfrak g\) be as in \SGARefLink{sga3:xiv:3.5}{3.5}. Then the
conditions \((\mathrm C_2)\), \((\mathrm C'_1)\), and
\((\mathrm C_1)\) of \SGARefLink{sga3:xiv:2.9}{2.9} are equivalent. In other words, if
\(\mathfrak g\) admits a Cartan subalgebra fpqc locally, then every
quasi-regular section of \(\mathfrak g\) is regular.}

Indeed, let \(a\) be a quasi-regular section; we prove that it is
regular. Since the question is fpqc local, we may suppose that
\(\mathfrak g\) admits a Cartan subalgebra \(\mathfrak d\). By
\SGARefLink{sga3:xiv:3.4}{3.4}\textup{(a)}, étale locally \(a\) is conjugate to a section of
\(\mathfrak d\). We are thereby reduced to the case in which
\(a\) is a section of \(\mathfrak d\), where the conclusion follows
immediately from the definition.

\paragraph{\textbf{Definition 3.8.}}
\SGARefTarget{sga3:xiv:3.8}
{\itshape Let \(G\) be a smooth prescheme in groups over a
prescheme \(S\). A \emph{subgroup of type \((\mathrm C)\)} of
\(G\) is a subprescheme in groups \(D\subset G\), smooth over
\(S\) and with connected fibers, such that
\[
  \mathfrak d=\operatorname{Lie}(D)
\]
is a Cartan subalgebra \textup{(\SGARefLink{sga3:xiv:2.4}{2.4})} of
\(\mathfrak g=\operatorname{Lie}(G)\). Equivalently, for every
\(s\in S\), the fiber \(D_s\) is a subgroup of type
\((\mathrm C)\) of the algebraic group \(G_s\)
\textup{(\SGARefLink{sga3:xiii:definition:6:2}{XIII~6.2})}.}

% END inputs/XIV/tex/components/11_expose_XIV_corollary34_through_definition38_en.tex

% BEGIN inputs/XIV/tex/components/12_expose_XIV_theorem39_through_corollary311_en.tex
% Authority: Expo14.pdf, local pp. 13--14 / combined pp. 847--848.
% Continuous English ends immediately before Proposition 3.12.

\paragraph{\textbf{Theorem 3.9.}}
\SGARefTarget{sga3:xiv:3.9}
{\itshape Let \(G\) be a smooth \(S\)-prescheme in groups of finite
presentation, and let \(\mathfrak g\) be its Lie algebra. Then:

\smallskip
\noindent
\textup{a)}
The map
\[
  D\longmapsto\mathfrak d=\operatorname{Lie}(D)
\]
establishes a bijective correspondence between the subgroups of
type \((\mathrm C)\) of \(G\) and the Cartan subalgebras of
\(\mathfrak g\).

\smallskip
\noindent
\textup{b)}
If \(D\) and \(\mathfrak d\) correspond, then
\[
  \operatorname{Norm}_G(D)=\operatorname{Norm}_G(\mathfrak d).
\]
This is a closed subprescheme of \(G\), smooth over \(S\), and
\[
  D=\operatorname{Norm}_G(D)^0
   =\operatorname{Norm}_G(\mathfrak d)^0.
\]

\smallskip
\noindent
\textup{c)}
Any two subgroups \(D\) and \(D'\) of type \((\mathrm C)\) in
\(G\) are conjugate étale locally.}

\emph{Proof.}
Let \(D\) be a subgroup of type \((\mathrm C)\) of \(G\), and let
\(\mathfrak d\) be its Lie algebra. Then
\[
  D\subset\operatorname{Norm}_G(\mathfrak d).
\]
By \SGARefLink{sga3:xiv:3.8}{Definition 3.8} and \SGARefLink{sga3:xiv:3.6}{Corollary 3.6}, this is an inclusion of smooth
preschemes in groups over \(S\) that induces an isomorphism on Lie
algebras. Since \(D\) has connected fibers, it follows that
\[
  D=\operatorname{Norm}_G(\mathfrak d)^0.
\]
Thus the map in \textup{a)} is injective. We prove that it is
surjective. Let \(\mathfrak d\) be a Cartan subalgebra of
\(\mathfrak g\). By \SGARefLink{sga3:xiv:3.6}{3.6},
\[
  N=\operatorname{Norm}_G(\mathfrak d)
\]
is a closed subprescheme in groups of \(G\), smooth over \(S\), and
has Lie algebra \(\mathfrak d\). Since \(G\) is of finite
presentation over \(S\), so is \(N\). Consequently, as noted in
XII after \SGARefLink{sga3:xii:7.3}{7.3}, the union of the connected components of the fibers
of \(N\) is the underlying set of an open subprescheme in groups
\(N^\circ\subset N\). Evidently \(N^\circ\) is a subgroup of type
\((\mathrm C)\) of \(G\) with Lie algebra \(\mathfrak d\). This
proves \textup{a)}. The first assertion of \textup{b)} follows at
once, and the formula
\[
  D=\operatorname{Norm}_G(\mathfrak d)^0
\]
has already been proved. Finally, \textup{c)} follows from
\textup{a)} and \SGARefLink{sga3:xiv:3.5}{3.5}.

\paragraph{\textbf{Corollary 3.10.}}
\SGARefTarget{sga3:xiv:3.10}
{\itshape Suppose that \(\mathfrak g\) admits a Cartan subalgebra
fpqc locally, or equivalently by \SGARefLink{sga3:xiv:3.9}{3.9}\textup{(a)}, that \(G\) admits
a subgroup of type \((\mathrm C)\) fpqc locally. Consider the
functor
\[
  \mathcal D\colon(\mathrm{Sch})^\circ/S\longrightarrow
  (\mathrm{Sets}),
\]
where \(\mathcal D(S')\) is the set of subgroups of type
\((\mathrm C)\) of \(G_{S'}\). This functor is representable by a
quasi-projective prescheme, smooth over \(S\), with connected
geometric fibers. If \(S=\operatorname{Spec}(k)\), so that \(G\)
is a smooth algebraic group over \(k\), and if \(\mathfrak g\)
admits a regular point---a condition automatically satisfied when
\(k\) is infinite---then \(\mathcal D\) is a rational variety over
\(k\).}

Indeed, by \SGARefLink{sga3:xiv:3.9}{3.9}\textup{(a)} the functor \(\mathcal D\) is
canonically isomorphic to the functor considered in \SGARefLink{sga3:xiv:2.16}{2.16}. On the
other hand, by \SGARefLink{sga3:xiv:3.7}{3.7} condition \((\mathrm C_2)\) is satisfied.
Therefore \SGARefLink{sga3:xiv:3.10}{3.10} follows from \SGARefLink{sga3:xiv:2.16}{2.16} and \SGARefLink{sga3:xiv:2.17}{2.17}.

\paragraph{\textbf{Corollary 3.11.}}
\SGARefTarget{sga3:xiv:3.11}
{\itshape Let \(D\) be a subgroup of type \((\mathrm C)\) of
\(G\), and let \(N\) be its normalizer in \(G\). Then the quotient
sheaf \(G/N\) is canonically isomorphic to the functor
\(\mathcal D\) of \SGARefLink{sga3:xiv:3.10}{3.10}. It is therefore representable by a
quasi-projective prescheme, smooth over \(S\), with connected
geometric fibers.}

% END inputs/XIV/tex/components/12_expose_XIV_theorem39_through_corollary311_en.tex

% BEGIN inputs/XIV/tex/components/13_expose_XIV_proposition312_through_corollary315_en.tex
% Authority: Expo14.pdf, local pp. 14--15 / combined pp. 848--849.
% Continuous English ends immediately before Corollary 3.16.

\paragraph{\textbf{Proposition 3.12.}}
\SGARefTarget{sga3:xiv:3.12}
{\itshape Let \(G\) be a smooth \(S\)-prescheme in groups of finite
presentation over \(S\). Let \(H\) and \(K\) be two subpreschemes
in groups, smooth and of finite presentation over \(S\), with
\(K\) having connected fibers, and let
\(\mathfrak g,\mathfrak h,\mathfrak k\) be the corresponding Lie
algebras. Suppose that one of the following conditions holds on
their geometric fibers:

\smallskip
\noindent
\textup{a)}
For every \(s\in S\), \(\mathfrak k_{\bar s}\) contains a regular
element of \(\mathfrak g_{\bar s}\), and
\(\mathfrak h_{\bar s}\) contains a Cartan subalgebra of
\(\mathfrak g_{\bar s}\).

\smallskip
\noindent
\textup{b)}
For every \(s\in S\), \(\mathfrak k_{\bar s}\) contains a Cartan
subalgebra of \(\mathfrak g_{\bar s}\).

\smallskip
\noindent
Under these conditions, \(H\supset K\) if and only if
\(\mathfrak h\supset\mathfrak k\).}

It remains only to prove that
\(\mathfrak h\supset\mathfrak k\) implies \(H\supset K\). In case
\textup{b)}, the inclusion
\(\mathfrak h\supset\mathfrak k\) places us under the hypotheses of
case \textup{a)}, so it is enough to prove \textup{a)}. Proceeding
as in \SGARefLink{sga3:xiv:3.2}{3.2}\textup{(a)}, by reduction to the case where \(S\) is
local artinian, we reduce to the case in which there exists a
section \(a\) of \(\mathfrak k\) that is quasi-regular in
\(\mathfrak g\). Arguing then as in \SGARefLink{sga3:xiii:corollary:5:5}{XIII~5.5} reduces us to the
following statement.

\paragraph{\textbf{Corollary 3.13.}}
\SGARefTarget{sga3:xiv:3.13}
{\itshape Let \(G\) be a smooth \(S\)-prescheme in groups of finite
presentation over \(S\), let \(H\subset G\) be a subprescheme in
groups smooth and of finite presentation over \(S\), and let
\(\mathfrak g\) and \(\mathfrak h\) be their Lie algebras. Let
\(a\) be a section of \(\mathfrak h\) that is quasi-regular in
\(\mathfrak g\). Suppose that for every \(s\in S\), the geometric
fiber \(\mathfrak h_{\bar s}\) contains a Cartan subalgebra of
\(\mathfrak g_{\bar s}\). Put
\[
  M_a=\operatorname{Transp}_G(a,\mathfrak h).
\]
This is a closed subprescheme of \(G\), smooth over \(S\)
\textup{(cf.\ \SGARefLink{sga3:xiv:3.1}{3.1})}. Hence \(M_a^0\), the union in the fibers of
\(M_a\) of the connected components containing the identity, is
an open subset of \(M_a\), endowed with the structure induced from
\(M_a\). Then
\[
  H^0=M_a^0.
\]}

Clearly \(H\subset M_a\), and therefore
\(H^0\subset M_a^0\). Since this is an inclusion of preschemes
smooth over \(S\), it is enough to check equality on the fibers.
This reduces us to the case where \(S\) is the spectrum of an
algebraically closed field, which was treated in \SGARefLink{sga3:xiii:corollary:5:4}{XIII~5.4}.

\paragraph{\textbf{Corollary 3.14.}}
\SGARefTarget{sga3:xiv:3.14}
{\itshape Let \(G\) be a smooth \(S\)-prescheme in groups of finite
presentation, let \(D\) be a subgroup of type \((\mathrm C)\) of
\(G\), and let \(H\subset G\) be a subprescheme in groups smooth
and of finite presentation over \(S\). If \(\mathfrak d\) and
\(\mathfrak h\) are their Lie algebras, then
\[
  \operatorname{Transp}_G(D,H)
  =
  \operatorname{Transp}_G(\mathfrak d,\mathfrak h),
\]
and this functor is representable by a closed subprescheme of
\(G\), smooth over \(S\).}

Indeed, the identity between the two transporters follows from
\SGARefLink{sga3:xiv:3.12}{3.12}\textup{(b)}, after which \SGARefLink{sga3:xiv:3.2}{3.2} applies.

\paragraph{\textbf{Corollary 3.15.}}
\SGARefTarget{sga3:xiv:3.15}
{\itshape Let \(G\) and \(H\) be as in \SGARefLink{sga3:xiv:3.14}{3.14}, and suppose that for
every \(s\in S\), the geometric fiber \(\mathfrak h_{\bar s}\)
contains a Cartan subalgebra of \(\mathfrak g_{\bar s}\). Suppose
also that \(\mathfrak g\) admits a Cartan subalgebra fpqc locally.
Then, étale locally, \(H\) contains a subgroup of type
\((\mathrm C)\) of \(G\).}

By \SGARefLink{sga3:xiv:3.7}{3.7} and \SGARefLink{sga3:xiv:3.9}{3.9}\textup{(a)}, \(G\) admits a subgroup of type
\((\mathrm C)\) étale locally. We may therefore suppose that \(G\)
admits such a subgroup \(D\). The hypothesis on \(\mathfrak h\)
then also says that the structure morphism of the transporter in
\SGARefLink{sga3:xiv:3.14}{3.14} is surjective, by the conjugacy theorem \SGARefLink{sga3:xiii:theorem:6:1}{XIII~6.1}\textup{(a)}.
The conclusion follows from Hensel's lemma \SGARefLink{sga3:xi:proposition:1:10}{XI~1.10}.

% END inputs/XIV/tex/components/13_expose_XIV_proposition312_through_corollary315_en.tex

% BEGIN inputs/XIV/tex/components/14_expose_XIV_corollary316_remarks317_en.tex
% Authority: Expo14.pdf, local pp. 15--16 / combined pp. 849--850.
% Continuous English ends immediately before Theorem 3.18.

\paragraph{\textbf{Corollary 3.16.}}
\SGARefTarget{sga3:xiv:3.16}
{\itshape Let \(G,H,K\) be as in \SGARefLink{sga3:xiv:3.12}{3.12}\textup{(a)}, and suppose in
addition that every geometric fiber \(\mathfrak k_{\bar s}\) is
nilpotent, that is, that \(\mathfrak k\) is locally nilpotent. Then
\[
  \operatorname{Transp}_G(K,H)
  =
  \operatorname{Transp}_G(\mathfrak k,\mathfrak h),
\]
and this functor is representable by a closed subprescheme of
\(G\), smooth over \(S\), whose structure morphism to \(S\) is
surjective. Étale locally, \(H\) contains a subgroup conjugate to
\(K\).}

The identity between the two transporters is again contained in
\SGARefLink{sga3:xiv:3.12}{3.12}\textup{(a)}. The assertion about its structure is precisely
\SGARefLink{sga3:xiv:3.2}{3.2}\textup{(a)}, and the last assertion of \SGARefLink{sga3:xiv:3.16}{3.16} then follows from
Hensel's lemma.

\paragraph{\textbf{Remarks 3.17.}}
\SGARefTarget{sga3:xiv:3.17}
\textup{a)}
In \SGARefLink{sga3:xiv:3.12}{3.12} and \SGARefLink{sga3:xiv:3.16}{3.16}, the hypothesis that \(K\) be smooth over \(S\)
may be replaced by the weaker hypothesis that the sheaf
\[
  e_K^*(\Omega^1_{K/S})=\omega^1_K
\]
of relative \(1\)-differentials along the identity section be
locally free. In this form, \SGARefLink{sga3:xiv:3.12}{3.12} includes \SGARefLink{sga3:xiii:corollary:5:5}{XIII~5.5}.

\smallskip
\noindent
\textup{b)}
Let \(G,\mathfrak g,\mathfrak h\) be as in \SGARefLink{sga3:xiv:3.1}{3.1}, with \(G\) of
finite presentation over \(S\). Then
\[
  N=\operatorname{Norm}_G(\mathfrak h)
\]
need not be smooth over \(S\) along the identity section.
Equivalently, there need not exist a subprescheme in groups
\(H\subset G\), smooth over \(S\), whose Lie algebra is
\(\mathfrak h\), even when \(S\) is the spectrum of a field. When
such an \(H\) exists, taking \(H\) with connected fibers gives
\[
  N=\operatorname{Norm}_G(H);
\]
I do not know whether \(N\) is then smooth over \(S\). In this
question one may plainly reduce to the case where \(S\) is local
artinian.

% END inputs/XIV/tex/components/14_expose_XIV_corollary316_remarks317_en.tex

% BEGIN inputs/XIV/tex/components/15_expose_XIV_theorem318_corollary319_en.tex
% Authority: Expo14.pdf, local pp. 16--17 / combined pp. 850--851.
% Continuous English ends immediately before Corollary 3.20.

To conclude this section, we examine the case in which \(G\) is
``semisimple'' over \(S\).

\paragraph{\textbf{Theorem 3.18.}}
\SGARefTarget{sga3:xiv:3.18}
{\itshape Let \(S\) be a prescheme and let \(G\) be a smooth
\(S\)-prescheme in groups whose geometric fibers are semisimple
``adjoint'' groups, that is, semisimple groups whose reductive
center \textup{(\SGARefLink{sga3:xii:4.1}{XII~4.1} and \SGARefLink{sga3:xiv:4.4}{4.4})} is the identity group. Then the
subgroups of type \((\mathrm C)\) of \(G\) are precisely its
maximal tori, and hence also precisely its Cartan subgroups
\textup{(\SGARefLink{sga3:xii:3.1}{XII~3.1})}.}

In view of the definitions, it is enough to treat the case where
\(S\) is the spectrum of an algebraically closed field. We must
then show that, for a maximal torus \(T\) of \(G\), its Lie algebra
\(\mathfrak t\) is a Cartan subalgebra of \(\mathfrak g\). Put
\(r=\dim T=\operatorname{rank}\mathfrak t\). Taking into account
the inequality
\[
  \operatorname{nil\,rank}(\mathfrak g)
  \geq
  \operatorname{nil\,rank}(G)
  =
  \dim T
  =
  \operatorname{rank}\mathfrak t
  =
  r,
\]
it is enough to find \(x\in\mathfrak t\) such that
\[
  \dim\Nil(x,\mathfrak g)=r.
\]
Since \(\mathfrak t\) is abelian, and therefore nilpotent, this is
equivalent to finding \(a\in\mathfrak t\) for which
\[
  \operatorname{ad}(a)_{\mathfrak g/\mathfrak t}
\]
is injective \textup{(\SGARefLink{sga3:xiii:corollary:5:7}{XIII~5.7}\textup{(a)})}.

Consider the characters \(\alpha\) of \(T\) occurring in the
representation of \(T\) induced by the adjoint representation of
\(G\). The structure theory of the semisimple group \(G\)
\textup{\emph{(Bible)}, 13, \SGARefLink{sga3:xiv:7:lemma1}{Theorem~1}\textup{(a)}, and Theorem~3,
\SGARefLink{sga3:xiv:7:lemma2}{Corollary~2}}, or more precisely the ``big cell'' of \(G\), shows
that the eigenspace of \(\mathfrak g\) for the trivial character is
exactly \(\mathfrak t\). The other eigenspaces are
one-dimensional, and their characters are precisely the roots of
\(G\) relative to \(T\). Here the big cell is built from \(T\) and
the subgroups \(P_\alpha\), each isomorphic to the additive group
\(\mathbf G_a\), stable under \(T\), and corresponding to the root
characters.

The reductive center of \(G\) is the intersection of the kernels
of the characters of \(T\) occurring in the adjoint representation
\textup{(\SGARefLink{sga3:xii:4.8}{XII~4.8})}. Thus the assertion that \(G\) is adjoint says
that its roots generate the lattice
\[
  M=\operatorname{Hom}(T,\mathbf G_m).
\]
A familiar lemma from root theory says that every root belongs to
a system of simple roots, hence to a basis of the group generated
by the roots, and therefore to a basis of the dual lattice of
\(T\).\footnote{The re-edition replaces \(\widehat T\) by \(T\)
at this point.} We obtain the following.

\paragraph{\textbf{Corollary 3.19.}}
\SGARefTarget{sga3:xiv:3.19}
{\itshape Let \(G\) be an adjoint semisimple algebraic group over
an algebraically closed field \(k\), and let \(T\) be a maximal
torus of \(G\). For every root
\[
  \alpha\colon T\longrightarrow\mathbf G_m
\]
of \(G\) relative to \(T\), the corresponding homomorphism
\[
  \alpha'\colon\mathfrak t\longrightarrow k
\]
is nonzero.}

This assertion is essentially equivalent to \SGARefLink{sga3:xiv:3.18}{Theorem 3.18}. Indeed,
for \(t\in\mathfrak t\), the endomorphism
\(\operatorname{ad}(t)\) is semisimple, and its eigenvalues on
\(\mathfrak g/\mathfrak t\) are precisely the numbers
\(\alpha'(t)\). Hence
\(\operatorname{ad}(t)_{\mathfrak g/\mathfrak t}\) is injective if
and only if all the \(\alpha'(t)\) are nonzero; and such a
\(t\in\mathfrak t\) exists if and only if every \(\alpha'\) is
nonzero.

% END inputs/XIV/tex/components/15_expose_XIV_theorem318_corollary319_en.tex

% BEGIN inputs/XIV/tex/components/16_expose_XIV_corollary320_remarks321_en.tex
% Authority: Expo14.pdf, local pp. 17--19 / combined pp. 851--853.
% This closes Section 3; continuous English next begins at Section 4.

\paragraph{\textbf{Corollary 3.20.}}
\SGARefTarget{sga3:xiv:3.20}
{\itshape Let \(S\) be a prescheme and let \(G\) be a smooth
\(S\)-prescheme in groups of finite presentation over \(S\), whose
geometric fibers are connected reductive algebraic groups, that is,
extensions of a semisimple group by a torus. For every \(s\in S\)
there is an open neighborhood \(U\) of \(s\) such that
\(G|_U\) admits a maximal torus.}\footnote{This result, as well as
\SGARefLink{sga3:xiv:2.11}{2.11} on which it rests, immediately generalizes to the case in
which \(s\) is replaced by a finite subset of \(S\) contained in
an affine open.}

We shall see in Exposé XVI that the hypothesis on \(G\) implies
that \(G\) is affine over \(S\) and has locally constant reductive
rank. Hence, by \SGARefLink{sga3:xii:4.7}{XII~4.7}\textup{\SGARefLink{sga3:xii:4.7:item:c}{(c)}}, \(G\) admits a reductive center
\(Z\), the quotient
\[
  G'=G/Z
\]
is a smooth affine group over \(S\) whose reductive center is the
identity group, and the maximal tori of \(G\) correspond
bijectively to those of \(G'\). Moreover, the geometric fibers of
\(G'\) are plainly connected semisimple groups, and are adjoint by
definition because their reductive center is trivial.

It is therefore enough to treat the case where \(G\) is
semisimple and adjoint. By \SGARefLink{sga3:xiv:3.18}{3.18}, the problem reduces to finding an
open neighborhood \(U\) of \(s\) and a subgroup of type
\((\mathrm C)\) of \(G|_U\), or equivalently by \SGARefLink{sga3:xiv:3.9}{3.9}\textup{(a)}, a
Cartan subalgebra of \(\mathfrak g|_U\). This is possible when
\(k(s)\) is infinite: by \SGARefLink{sga3:xiv:3.7}{3.7}, \(\mathfrak g\) satisfies condition
\((\mathrm C_2)\) of \SGARefLink{sga3:xiv:2.9}{2.9}, and \SGARefLink{sga3:xiv:2.11}{2.11}\textup{(b)} applies.

In fact, \SGARefLink{sga3:xiv:3.20}{3.20} remains true without assuming that \(k(s)\) is
infinite. By the preceding argument, it is enough to know that, for
every adjoint semisimple group \(G\) over a finite field \(k\), the
Lie algebra \(\mathfrak g\) contains a regular element. Chevalley
proved this assertion using the properties of a Coxeter element of
the Weyl group; see the appendix below by J.-P. Serre.

\paragraph{\textbf{Remarks 3.21.}}
\SGARefTarget{sga3:xiv:3.21}
\textup{a)}
Assertion \SGARefLink{sga3:xiv:3.20}{3.20} remains valid, with essentially the same proof, if
\(G\) is replaced by a closed subprescheme in groups \(H\), smooth
over \(S\), whose rank is everywhere the same as that of \(G\)
\textup{(for example, a ``parabolic subgroup'' of \(G\))}, provided
that \(k(s)\) is infinite. I do not know whether the infinitude
hypothesis is again superfluous. It can at least be removed when
\(H\) is parabolic: constructing the radical \(U\) of \(H\) and
the quotient \(H/U\) reduces the question to the semisimple case.
Unfortunately, the method of regular elements seems powerless
here. One readily constructs examples---for instance, the
projective group with its ``standard'' Borel subgroup---in which
the Lie algebra \(\mathfrak h\) of \(H\) contains no regular
element.

\smallskip
\noindent
\textup{b)}
When \(k(s)\) is infinite, the proof of \SGARefLink{sga3:xiv:3.20}{3.20} in fact gives the
following sharper statement, using \SGARefLink{sga3:xiv:3.9}{3.9}\textup{(b)}: every maximal
torus \(T_0\) of \(G_s\) is induced by a maximal torus \(T\) over
an open neighborhood of \(s\). I do not know whether this remains
true without the infinitude hypothesis. The evident difficulty is
that the Lie algebra \(\mathfrak t_0\) of \(T_0\) need not contain
a regular element of the Lie algebra \(\mathfrak g_0\) of the
fiber \(G_s\).

An affirmative answer would imply the following assertion, which
is known only when the residue field is infinite or when \(A\) is
separated and complete. Let \(A\) be a local ring with residue
field \(k\), and let \(M\) be an ``Azumaya algebra'' over \(A\),
meaning an algebra that is a finite free \(A\)-module and whose
reduction
\[
  M_0=M\otimes_A k
\]
is a central simple algebra over \(k\). Let \(D_0\subset M_0\) be
a commutative subalgebra separable over \(k\), with
\[
  [M_0:k]=[D_0:k]^2.
\]
Then there should exist a commutative subalgebra \(D\subset M\)
that is a direct summand as an \(A\)-module and satisfies
\[
  D\otimes_A k=D_0.
\]

The datum of \(M\) is equivalent to that of a principal homogeneous
bundle under the projective group \(\operatorname{PGL}(n)_A\), and
hence to an inner twisted form \(G\) of \(\operatorname{PGL}(n)\).
The maximal tori of \(G\) correspond bijectively to the commutative
subalgebras \(D\subset M\) that are étale of rank \(n\) over \(A\).

\smallskip
\noindent
\textup{c)}
Applying \SGARefLink{sga3:xiv:3.20}{3.20} to the centralizer of a subtorus \(Q\) of a reductive
\(S\)-prescheme in groups \(G\), one sees that every such \(Q\) is
contained, Zariski locally, in a maximal torus of \(G\).

% END inputs/XIV/tex/components/16_expose_XIV_corollary320_remarks321_en.tex

% BEGIN inputs/XIV/tex/components/17_expose_XIV_section4_definition41_through_definition45_en.tex
% Authority: Expo14.pdf, local pp. 18--19 / combined pp. 852--853.
% Exact next cursor: Remarks 4.6 on local p. 19 / combined p. 853.

\section{A Digression on Borel Subgroups}
\SGARefTarget{sga3:xiv:4}

\paragraph{\textbf{Definition 4.1.}}
\SGARefTarget{sga3:xiv:4.1}
{\itshape Let \(G\) be a smooth algebraic group over an algebraically
closed field. A \emph{Borel subgroup} of \(G\) is a smooth, solvable,
connected algebraic subgroup that is maximal with these properties.}

When \(G\) is affine, this agrees with the terminology of
\textup{\emph{(Bible)}, 6, \SGARefLink{sga3:xiv:7:lemma1}{Definition~1}}. Observe at once that if
\(Z\) is a connected smooth subgroup of \(G\) contained in its center
\textup{(or, more generally, if \(Z\) is solvable and normal)}, then,
for every Borel subgroup \(B\) of \(G\), the image \(BZ\) of
\(B\times Z\) under the morphism
\[
  (b,z)\longmapsto bz
\]
from \(B\times Z\) to \(G\) is a smooth, solvable, connected subgroup
of \(G\) containing \(B\). It is therefore equal to \(B\). Thus \(B\)
contains \(Z\), and consequently \(B\) is the inverse image of an
algebraic subgroup \(B'\) of \(G'=G/Z\); it is immediate that \(B'\)
is a Borel subgroup of \(G'\).

Taking
\[
  Z=\Centr(G^0)_{\mathrm{red}}^0,
\]
the quotient \(G'=G/Z\) is affine by \SGARefLink{sga3:xii:6.1}{XII~6.1}. Hence the Borel
subgroups of \(G'\) are conjugate, and \(G'/B'\) is a projective
variety for every such \(B'\)
\textup{\emph{(Bible)}, 6, \SGARefLink{sga3:xiv:7:lemma4}{Theorem~4}}. Consequently:

\paragraph{\textbf{Proposition 4.2.}}
\SGARefTarget{sga3:xiv:4.2}
{\itshape Let \(G\) be as in \SGARefLink{sga3:xiv:4.1}{4.1}. Then the Borel subgroups of \(G\)
are conjugate. If \(B\) is a Borel subgroup, then \(G/B\) is a
projective variety. The maximal tori of \(B\) are maximal tori of
\(G\); and, when \(G\) is connected, the Cartan subgroups of \(B\)
are Cartan subgroups of \(G\).}

It remains to prove the last assertion, and we may plainly suppose
that \(G=G^0\). For Cartan subgroups, it follows from the analogous
assertion in \(G'\)
\textup{\emph{(Bible)}, 6, \SGARefLink{sga3:xiv:7:lemma4}{Theorem~4}, \SGARefLink{sga3:xiv:7:lemma4}{Corollary~4}}
and \SGARefLink{sga3:xii:6.6}{XII~6.6}\textup{\SGARefLink{sga3:xii:6.6:item:e}{(e)}}. The assertion for maximal tori follows from
this one, since by \SGARefLink{sga3:xii:6.6}{XII~6.6}\textup{\SGARefLink{sga3:xii:6.6:item:c}{(c)}} the maximal tori of a smooth
algebraic group are precisely the maximal tori of its Cartan
subgroups.

\paragraph{\textbf{Corollary 4.3.}}
\SGARefTarget{sga3:xiv:4.3}
{\itshape Suppose that \(G\) is connected. Then every element of
\(G\) is contained in a Borel subgroup of \(G\).}

This reduces to the same assertion in \(G'\), where it is well known
\textup{\emph{(Bible)}, 6, Theorem~5\textup{(d)}}.

\paragraph{\textbf{Corollary 4.4.}}
\SGARefTarget{sga3:xiv:4.4}
{\itshape Let \(B\) be a Borel subgroup of \(G\), let \(C\) be a
Cartan subgroup of \(B\), and let \(N\) be its normalizer in \(G\).
Then
\[
  N\cap B=C.
\]}

Indeed,
\[
  N\cap B=\operatorname{Norm}_B(C),
\]
so it remains to show that, when \(G\) is connected and solvable, a
Cartan subgroup \(C\) is its own connected normalizer. With the
preceding notation, \(C\) is the inverse image of a Cartan subgroup
\(C'\) of \(G'\), and we are therefore reduced to the case in which
\(G\) is affine. The normalizer of a Cartan subgroup is smooth
\textup{(\(C\) being its own connected normalizer; see, for example,
\SGARefLink{sga3:xii:6.6}{XII~6.6}\textup{\SGARefLink{sga3:xii:6.6:item:c}{(c)}})}. It is thus enough to check that \(C\) and
\(N\) have the same \(k\)-valued points, which is exactly
\textup{\emph{(Bible)}, Theorem~6\textup{\SGARefLink{sga3:xii:6.6:item:d}{(d)}}}.

\paragraph{\textbf{Definition 4.5.}}
\SGARefTarget{sga3:xiv:4.5}
{\itshape Let \(G\) be a smooth prescheme in groups of finite
presentation over a prescheme \(S\). A \emph{Borel subgroup} of
\(G\) is a smooth subprescheme in groups \(B\) of \(G\), of finite
presentation, such that for every \(s\in S\) the geometric fiber
\(B_{\bar s}\) is a Borel subgroup of \(G_{\bar s}\).}

As is immediately verified, this notion is stable under base change
and is local for the fpqc topology. Indeed, if \(k'\) is an
algebraically closed extension of an algebraically closed field
\(k\), then an algebraic subgroup \(B\) of a smooth algebraic group
\(G\) over \(k\) is a Borel subgroup of \(G\) if and only if
\(B_{k'}\) is a Borel subgroup of \(G_{k'}\).

It follows from the definition that, if \(G\) is a smooth algebraic
group over an arbitrary field \(k\) and \(B\) is a Borel subgroup
of \(G\), then \(G/B\) is a projective variety; every maximal torus
\(T\) of \(B\) is a maximal torus of \(G\); and the normalizer of
\(T\) in \(B\) equals its centralizer \(C\), which is a Cartan
subgroup of \(G\) when \(G\) is connected.

% END inputs/XIV/tex/components/17_expose_XIV_section4_definition41_through_definition45_en.tex

% BEGIN inputs/XIV/tex/components/18_expose_XIV_remarks46_through_remarks481_en.tex
% Authority: Expo14.pdf, local pp. 19--20 / combined pp. 853--854.
% The source has no item numbered 4.7.
% Exact next cursor: the paragraph preceding Proposition 4.9,
% local p. 20 / combined p. 854.

\paragraph{\textbf{Remarks 4.6.}}
\SGARefTarget{sga3:xiv:4.6}
Unfortunately, it is no longer true in general that two Borel
subgroups of \(G\) are locally conjugate for the fpqc topology.
This can fail even when \(G\) is affine over \(S\) and \(S\) is
the spectrum of the algebra of dual numbers over an algebraically
closed field \(k\).

As a consequence of this regrettable fact, let us point out that
if \(G\) is a smooth, affine, connected algebraic group over a
nonperfect field \(k\), it is in general impossible to define
naturally a homogeneous space \(D\) under \(G\) that plays the
role of a flag variety, that is, of the variety of Borel subgroups
of \(G\). Over the algebraic closure \(\bar k\) of \(k\), such a
space would have to be isomorphic to \(G_{\bar k}/B\), where \(B\)
is a Borel subgroup of \(G_{\bar k}\).

Indeed, suppose that the quotient of \(G_{\bar k}\) by its radical
\(\bar R\) is an adjoint semisimple group. The kernel of
\[
  G_{\bar k}\longrightarrow
  \operatorname{Aut}_{\bar k}(\bar D)
\]
is then the radical of \(G_{\bar k}\). Consequently, if \(\bar D\)
came from a homogeneous space \(D\) under \(G\), the radical
\(\bar R\) would come from a subgroup \(R\) of \(G\). But examples
are easily constructed for which \(G_{\bar k}/\bar R\) is adjoint
while \(\bar R\) is not defined over \(k\).

Under these conditions, the functor
\[
  \mathfrak B\colon(\operatorname{Sch}/S)^\circ
    \longrightarrow \operatorname{Ens},
  \qquad
  \mathfrak B(S')=
  \{\text{Borel subgroups of }G_{S'}\},
\]
is not representable by a smooth \(S\)-prescheme. From the
infinitesimal point of view \textup{(III, \S3)}, the failure of the
conjugacy theorem is expressed by the fact that, if \(B\) is a
Borel subgroup of the smooth algebraic group \(G\), the cohomology
group
\[
  H^1(B,\mathfrak g/\mathfrak b)
\]
may be nonzero.\footnote{Here \(\mathfrak b\) denotes the Lie
algebra of \(B\).}

We shall see, on the other hand, in a later exposé that when \(G\)
is semisimple, or more generally reductive, such unpleasant
phenomena do not occur. These phenomena, together no doubt with
the absence of good existence theorems, explain why Borel
subgroups play only a relatively unobtrusive role in the
scheme-theoretic study of general group schemes, whereas they
dominate the theory of semisimple group schemes in the later
exposés.

\paragraph{\textbf{Proposition 4.8.}}
\SGARefTarget{sga3:xiv:4.8}
\footnote{There is no item numbered 4.7 in the source.}
{\itshape Let \(G\) be a smooth \(S\)-prescheme in groups of finite
presentation with connected fibers, and let \(B\) be a Borel
subgroup of \(G\). Then \(B\) is its own normalizer and is a closed
subprescheme of \(G\).}

By \SGARefLink{sga3:xii:7.10}{XII~7.10}, it is enough to prove that, over an algebraically
closed field \(k\), every element of \(G(k)\) that normalizes \(B\)
belongs to \(B(k)\). When \(G\) is affine this is a fundamental
result of Chevalley
\textup{\emph{(Bible)}, 9, \SGARefLink{sga3:xiv:7:lemma1}{Theorem~1}}; the general case reduces
immediately to it by the reduction already used in \SGARefLink{sga3:xiv:4.2}{4.2}.

\paragraph{\textbf{Remarks 4.8.1.}}
\SGARefTarget{sga3:xiv:4.8.1}
\SGARefLink{sga3:xiv:4.5}{Definition~4.5} may be generalized by also introducing the notion
of a \emph{parabolic subgroup} of \(G\). This means a subprescheme
in groups \(P\) of \(G\), smooth and of finite presentation over
\(S\), such that for every \(s\in S\) the geometric fiber
\(P_{\bar s}\) is a parabolic subgroup of \(G_{\bar s}\), that is,
contains a Borel subgroup of \(G_{\bar s}\).

\SGARefLink{sga3:xiv:4.8}{Proposition~4.8} extends to a parabolic subgroup \(P\) of \(G\), by
the same reduction to the known set-theoretic assertion. We note
the following consequence \textup{(compare Exposé XVI)}. \SGAInlineRefTarget{sga3:xiv:paragraph:4-8-bis}If \(P\)
is a parabolic subgroup of \(G\), then \(G/P\) is representable by
a quasi-projective prescheme of finite presentation over \(S\);
here \(G\) is assumed to have connected fibers. Moreover, \(G/P\)
is plainly smooth over \(S\), and its geometric fibers are
connected and proper. It follows readily from EGA~III, 5.5.1 that
\[
  D=G/P
\]
is proper, hence projective, over \(S\). If its relative dimension
is \(n\), it is known that the inverse of the invertible sheaf
\(\Omega^n_{D/S}\) induces ample sheaves on the geometric fibers of
\(D/S\); hence, by EGA~III, 4.7.1,
\[
  \bigl(\Omega^n_{D/S}\bigr)^{-1}
\]
is ample on \(D\) relative to \(S\).

% END inputs/XIV/tex/components/18_expose_XIV_remarks46_through_remarks481_en.tex

% BEGIN inputs/XIV/tex/components/19_expose_XIV_proposition49_corollary410_en.tex
% Authority: Expo14.pdf, local pp. 20--21 / combined pp. 854--855.
% Exact next cursor: Theorem 4.11, local p. 21 / combined p. 855.

It is readily seen, by reduction to the affine case and to an
algebraically closed base field, that if
\[
  u\colon G\longrightarrow G'
\]
is an epimorphism of smooth algebraic groups, then \(u(B)=B'\) is a
Borel subgroup of \(G'\) for every Borel subgroup \(B\) of \(G\).
We now consider the case in which this gives a bijective
correspondence between the Borel subgroups of \(G\) and those of
\(G'\).

\paragraph{\textbf{Proposition 4.9.}}
\SGARefTarget{sga3:xiv:4.9}
{\itshape Let \(G\) and \(G'\) be smooth \(S\)-preschemes in groups
of finite presentation with connected fibers, and let
\[
  u\colon G\longrightarrow G'
\]
be a faithfully flat, that is, surjective, homomorphism of
groups.\footnote{The editors ask whether a reference should be
added here.} Put \(N=\ker u\), and suppose that one of the following
two conditions holds:

\smallskip
\noindent
\textup{a)} \(N\) is central in \(G\).

\smallskip
\noindent
\textup{b)} \(S=\operatorname{Spec}(k)\) for a field \(k\), and, if
\(\bar k\) is an algebraic closure of \(k\), then
\[
  N_{\bar k}=\ker(u_{\bar k})
\]
is contained in the radical of \(G_{\bar k}\), that is, in its
largest smooth, connected, solvable, normal subgroup.

\smallskip
Then the map
\[
  B'\longmapsto u^{-1}(B')
\]
induces a bijection from the set of Borel subgroups of \(G'\) onto
the set of Borel subgroups of \(G\).}

Case \textup{b)} follows immediately from the correspondence between
algebraic subgroups of \(G'\) and algebraic subgroups of \(G\)
containing \(N\), together with the fact that, over an algebraically
closed field, the Borel subgroups of \(G\) contain its radical. The
latter fact follows at once from the argument preceding \SGARefLink{sga3:xiv:4.2}{4.2}.

We prove case \textup{a)}. By \SGARefLink{sga3:xii:7.12}{XII~7.12}, the map
\[
  H'\longmapsto H=u^{-1}(H')
\]
gives a bijection between the subpreschemes in groups \(H'\) of
\(G'\) that are smooth and of finite presentation over \(S\), have
connected fibers, and at every \(s\in S\) have the same reductive
rank and the same nilpotent rank as \(G'\), and the subpreschemes
in groups \(H\) of \(G\) having the analogous properties. Borel
subgroups of \(G'\) and of \(G\) have these properties.

It remains to show that, when \(H'\) and \(H\) correspond, \(H'\)
is a Borel subgroup of \(G'\) if and only if \(H\) is a Borel
subgroup of \(G\). By definition, this question reduces to the case
where \(S\) is the spectrum of an algebraically closed field. Since
\(N\) is central in \(G\), and hence in \(H\), the group \(H\) is
solvable if and only if \(H'\) is solvable. Finally, the
correspondence between algebraic subgroups of \(G'\) and algebraic
subgroups of \(G\) containing \(N\) shows immediately that \(H'\)
has the maximality property of \SGARefLink{sga3:xiv:4.1}{Definition~4.1} if and only if \(H\)
does. This completes the proof.

\paragraph{\textbf{Corollary 4.10.}}
\SGARefTarget{sga3:xiv:4.10}
{\itshape With the notation of \SGARefLink{sga3:xiv:4.9}{4.9},\footnote{The source prints
``\SGARefLink{sga3:xiv:4.9}{4.7}''; the intended reference is \SGARefLink{sga3:xiv:4.9}{4.9}.} if \(B'\) and \(B\) are
corresponding Borel subgroups of \(G'\) and \(G\), then
\[
  \mathfrak b=\operatorname{Lie}(u)^{-1}(\mathfrak b'),
\]
where \(\mathfrak g,\mathfrak g',\mathfrak b,\mathfrak b'\) are
the Lie algebras of \(G,G',B,B'\), respectively, and
\[
  \operatorname{Lie}(u)\colon\mathfrak g\longrightarrow\mathfrak g'
\]
is the homomorphism induced by \(u\).}

This follows immediately from the definitions and the identity
\[
  u^{-1}(B')=B.
\]

We can now prove the principal result of this section.

% END inputs/XIV/tex/components/19_expose_XIV_proposition49_corollary410_en.tex

% BEGIN inputs/XIV/tex/components/20_expose_XIV_theorem411_corollary412_en.tex
% Authority: Expo14.pdf, local pp. 21--22 / combined pp. 855--856.
% This closes Section 4.
% Exact next cursor: Section 5 and Theorem 5.1,
% local p. 22 / combined p. 856.

\paragraph{\textbf{Theorem 4.11.}}
\SGARefTarget{sga3:xiv:4.11}
{\itshape Let \(G\) be a smooth algebraic group over an
algebraically closed field \(k\), and let \(\mathfrak g\) be its
Lie algebra. Then \(\mathfrak g\) is the union of the Lie algebras
\(\mathfrak b\) of the Borel subgroups \(B\) of \(G\).}

We may plainly suppose that \(G\) is connected. Let \(R\) be the
radical of \(G\), and put \(G'=G/R\). By \SGARefLink{sga3:xiv:4.9}{4.9}\textup{(b)} and \SGARefLink{sga3:xiv:4.10}{4.10},
it is enough to prove the theorem for \(G'\); thus we may suppose
that \(G\) is semisimple.

Let \(Z\) be the center of \(G\), which is its reductive center, and
put \(G'=G/Z\). The same argument, now using \SGARefLink{sga3:xiv:4.9}{4.9}\textup{(a)},
reduces the theorem to \(G'\). We may therefore suppose that \(G\)
is semisimple and adjoint.

Let \(B\) be a Borel subgroup of \(G\), let \(T\) be a maximal torus
of \(B\), hence of \(G\), and let \(\mathfrak b\) and
\(\mathfrak t\) be their Lie algebras. By \SGARefLink{sga3:xiv:3.18}{3.18}, \(T\) is a subgroup
of type \((\mathrm C)\) of \(G\), that is, \(\mathfrak t\) is a
Cartan subalgebra of \(\mathfrak g\). The union of the conjugates
of \(\mathfrak t\) is consequently dense in \(\mathfrak g\)
\textup{(\SGARefLink{sga3:xiii:theorem:5:1}{XIII~5.1}, \textup{(i)}\(\Rightarrow\)\textup{(vii)})}.
\emph{A fortiori}, the union of the conjugates of \(\mathfrak b\)
is dense in \(\mathfrak g\).

Let \(X\) be the closed subscheme of
\[
  G/B\times W(\mathfrak g)
\]
whose \(k\)-valued points are the pairs \((gB,x)\) such that
\[
  x\in\operatorname{Ad}(g)\mathfrak b
\]
\textup{(compare XIII~1)}. The morphism
\[
  \psi\colon X\longrightarrow W(\mathfrak g)
\]
induced by the second projection is proper because \(G/B\) is
proper over \(k\). We have just seen that it is dominant; it is
therefore surjective. This proves \SGARefLink{sga3:xiv:4.11}{4.11}.

The only result of this section that we shall use in the remainder
of this exposé is the following corollary.

\paragraph{\textbf{Corollary 4.12.}}
\SGARefTarget{sga3:xiv:4.12}
{\itshape Let \(k\) be an infinite field, let \(G\) be a smooth
algebraic group over \(k\), and let \(T\) be a maximal torus of
\(G\). Write \(\mathfrak g\) and \(\mathfrak t\) for their Lie
algebras. Let
\[
  u\colon G\longrightarrow\operatorname{GL}(V)
\]
be a linear representation, where \(V\) is finite-dimensional over
\(k\), and let
\[
  u'\colon\mathfrak g\longrightarrow\mathfrak{gl}(V)
\]
be the induced representation of Lie algebras. Then the minimum of
the nullity of \(u'(x)\), for \(x\in\mathfrak g\), is attained at
an element \(x\in\mathfrak t\).}

We reduce immediately to the case where \(k\) is algebraically
closed. We may plainly suppose that \(G\) is connected, and, after
dividing \(G\) by \((\ker u)_{\mathrm{red}}\), that \(G\) is
affine. Using \SGARefLink{sga3:xiv:4.11}{4.11} and the conjugacy theorem for the maximal tori
of \(G\), we reduce to the further case in which \(G\) is solvable.

Then \(G\) is a semidirect product \(T\cdot U\), where \(U\) is the
unipotent part of \(G\), a smooth connected unipotent group
\textup{\emph{(Bible)}, 6, Theorem~3}. Thus, as a vector space,
\[
  \mathfrak g=\mathfrak t\oplus\mathfrak u,
  \qquad
  \mathfrak u=\operatorname{Lie}(U).
\]
By the Lie--Kolchin theorem
\textup{\emph{(Bible)}, 6, \SGARefLink{sga3:xiv:7:lemma1}{Theorem~1}}, the representation space
\(V\) has a composition series by stable subspaces \(V_i\), with
one-dimensional quotients
\[
  W_i=V_i/V_{i+1}.
\]
For every \(i\) there is therefore an induced representation
\[
  u_i\colon G\longrightarrow
  \operatorname{GL}(W_i)=\mathbb G_m
\]
and a corresponding homomorphism of Lie algebras
\[
  u_i'\colon\mathfrak g\longrightarrow k.
\]
For every \(x\in\mathfrak g\), the nullity of \(u'(x)\) is the
number of indices \(i\) for which \(u_i'(x)=0\).

Since \(U\) is unipotent, every \(u_i\) is trivial on \(U\), and
hence every \(u_i'\) vanishes on \(\mathfrak u\). If
\[
  x=t+v,\qquad t\in\mathfrak t,\quad v\in\mathfrak u,
\]
then \(u_i'(x)=u_i'(t)\) for every \(i\). Thus \(u'(x)\) and
\(u'(t)\) have the same nullity, and \SGARefLink{sga3:xiv:4.12}{4.12} follows.

% END inputs/XIV/tex/components/20_expose_XIV_theorem411_corollary412_en.tex

% BEGIN inputs/XIV/tex/components/21_expose_XIV_section5_theorem51_through_lemma53_en.tex
% Authority: Expo14.pdf, local pp. 22--24 / combined pp. 856--858.
% Exact next cursor: Corollary 5.4, local p. 24 / combined p. 858.

\section{Relations Between Cartan Subgroups and Cartan Subalgebras}
\SGARefTarget{sga3:xiv:5}

Applying \SGARefLink{sga3:xiv:4.12}{4.12} to the adjoint representation of \(G\), we obtain:

\paragraph{\textbf{Theorem 5.1.}}
\SGARefTarget{sga3:xiv:5.1}
{\itshape Let \(G\) be a smooth algebraic group over an infinite
field, let \(T\) be a maximal torus of \(G\), and let
\(\mathfrak g\supset\mathfrak t\) be their Lie algebras. Then
\(\mathfrak t\) contains a regular element of \(\mathfrak g\).}

\paragraph{\textbf{Corollary 5.2.}}
\SGARefTarget{sga3:xiv:5.2}
{\itshape Let \(G\) be a smooth \(S\)-prescheme in groups of finite
presentation, and let \(\mathfrak g\) be its Lie algebra.

\smallskip
\noindent
\textup{a)} Conditions \((\mathrm C_0)\) through \((\mathrm C_2)\)
of \SGARefLink{sga3:xiv:2.9}{2.9} for \(\mathfrak g\) are equivalent. In particular, if the
infinitesimal rank of the fibers of \(\mathfrak g\) at the points
of \(S\) is locally constant, then, étale locally,
\(\mathfrak g\) admits a Cartan subalgebra and hence, by
\SGARefLink{sga3:xiv:3.9}{3.9}\textup{(a)}, \(G\) admits a subgroup of type \((\mathrm C)\).

\smallskip
\noindent
\textup{b)} Let \(H\) be a subprescheme in groups of \(G\), smooth
and of finite presentation over \(S\), whose fibers are connected
and at every \(s\in S\) have the same reductive rank as \(G\);
for example, \(H\) may be a maximal torus or a Cartan subgroup of
\(G\). Let \(D\) be a subgroup of type \((\mathrm C)\) of \(G\).
Then
\[
  H\subset D\quad\Longleftrightarrow\quad
  \mathfrak h\subset\mathfrak d.
\]

\smallskip
\noindent
\textup{c)} Suppose that condition \((\mathrm C_0)\) holds, that is,
the infinitesimal rank of \(G\) is locally constant. Let \(H\) be a
subprescheme in groups of \(G\), smooth and of finite presentation
over \(S\), with connected nilpotent fibers, and having at every
point the same reductive rank as \(G\); for example, \(H\) may be a
maximal torus or a Cartan subgroup of \(G\). Then \(H\) is contained
in a unique subgroup \(D\) of type \((\mathrm C)\) of \(G\).

\smallskip
\noindent
\textup{d)} If \(G\) admits a Cartan subgroup fpqc locally, then so
does every subgroup \(D\) of type \((\mathrm C)\) of \(G\).}

\smallskip
\noindent
\emph{Proof.}
\textup{a)}
Assume \((\mathrm C_0)\), and let us prove \((\mathrm C_2)\), that
is, that every quasi-regular section \(a\) of \(\mathfrak g\) is
regular. As usual, we reduce first to \(S\) affine noetherian and
then, since the question is infinitesimal by \SGARefLink{sga3:xiv:2.15}{2.15}\textup{(b)}, to
\(S\) local Artinian. In this last case \((\mathrm C_0)\) is
automatic. We may moreover suppose that the residue field \(k\)
of \(S\) is infinite.

By \SGARefLink{sga3:xiv:3.7}{3.7}, it is enough to show that \(\mathfrak g\) admits a Cartan
subalgebra. Let \(T\) be a maximal torus of \(G\). By \SGARefLink{sga3:xiv:5.1}{5.1} there is
a quasi-regular element \(t\) in the Lie algebra \(\mathfrak t\)
of \(T\). We shall show that \(t\) is regular.

Consider the linear representation of \(T\) on \(\mathfrak g\)
induced by the adjoint representation of \(G\). There is a finite
set of characters \((u_i)_{i\in I}\) of \(T\) and a direct-sum
decomposition
\[
  \mathfrak g=\bigoplus_{i\in I}\mathfrak g_i
\]
into \(T\)-stable submodules such that \(T\) acts on
\(\mathfrak g_i\) through \(u_i\)
\textup{(I, \S4.7.3)}. Let
\[
  u_i'\colon\mathfrak t\longrightarrow A
\]
be the homomorphism of Lie algebras induced by
\(u_i\colon T\to\mathbb G_m\), where \(A\) is the ring of \(S\).
Let \(u_{i0}\) and \(u_{i0}'\) be their reductions under
\(A\to k\). Put
\[
  I'=\{i\in I\mid u_{i0}'\ne0\},
  \qquad I''=I\setminus I'.
\]
The regularity of \(t_0\) is expressed by
\[
  u_{i0}'(t_0)\ne0\qquad(i\in I'),
\]
and therefore \(u_i'(t)\) is invertible for \(i\in I'\).

The Cartan subalgebra of \(\mathfrak g_0\) defined by \(t_0\), that
is, the kernel of the semisimple endomorphism
\(\operatorname{ad}(t_0)\), is
\[
  \bigoplus_{i\in I''}(\mathfrak g_i)_0.
\]
Set
\[
  \mathfrak d=\bigoplus_{i\in I''}\mathfrak g_i.
\]
Then \(\operatorname{ad}(t)\) is nilpotent on \(\mathfrak d\), and
is an automorphism of
\[
  \mathfrak g/\mathfrak d
  \simeq\bigoplus_{i\in I'}\mathfrak g_i.
\]
By \SGARefLink{sga3:xiv:2.6}{2.6}, \(t\) is regular.

\smallskip
\noindent
\textup{b)}
This follows from \SGARefLink{sga3:xiv:3.12}{3.12}\textup{(a)}. By \SGARefLink{sga3:xiv:5.1}{5.1}, \(H\) satisfies the
required hypothesis that every geometric fiber of \(\mathfrak h\)
contains a regular element of the corresponding fiber of
\(\mathfrak g\).

\smallskip
\noindent
\textup{c)}
By \textup{b)}, it is enough to show that \(\mathfrak h\) is
contained in a unique Cartan subalgebra of \(\mathfrak g\). By
\textup{a)}, \(\mathfrak g\) satisfies \((\mathrm C_2)\). We are
therefore reduced to the following lemma.

\paragraph{\textbf{Lemma 5.3.}}
\SGARefTarget{sga3:xiv:5.3}
{\itshape Let \(\mathfrak g\) be a Lie algebra over \(S\), locally
free of finite type as a module, and satisfying condition
\((\mathrm C_2)\) of \SGARefLink{sga3:xiv:2.9}{2.9}. Let \(\mathfrak h\) be a Lie subalgebra
of \(\mathfrak g\) such that \(\mathfrak h\) is locally a direct
summand, is locally nilpotent, and, for every \(s\in S\), the
geometric fiber \(\mathfrak h_{\bar s}\) contains a regular element
of \(\mathfrak g_{\bar s}\). Then \(\mathfrak h\) is contained in
a unique Cartan subalgebra of \(\mathfrak g\).}

In the situation at hand, \(\mathfrak h\) has these properties: it
is locally nilpotent because \(H\) has nilpotent fibers, and the
condition on regular elements follows from \SGARefLink{sga3:xiv:5.1}{5.1}.

Since \SGARefLink{sga3:xiv:5.3}{5.3} is fpqc local, it is enough, at a point \(s\in S\) with
\(\kappa(s)\) infinite, to find an open neighborhood \(U\) of \(s\)
on which existence and uniqueness hold after every base change
\(S'\to S\) factoring through \(U\). Choose a regular element of
\(\mathfrak g\otimes\kappa(s)\) lying in
\(\mathfrak h\otimes\kappa(s)\), and extend it to a section \(a\)
of \(\mathfrak h\) over an open neighborhood \(U\) of \(s\).
Condition \((\mathrm C_2)\) and \SGARefLink{sga3:xiv:2.10}{2.10} allow us to shrink \(U\) so
that \(a\) is regular. We may thus suppose \(U=S\).

Every Cartan subalgebra of \(\mathfrak g\) containing
\(\mathfrak h\) contains \(a\), and hence equals
\[
  \mathfrak d=\Nil(a,\mathfrak g)
\]
by \SGARefLink{sga3:xiv:2.6}{2.6}; this proves uniqueness. Conversely,
\(\mathfrak h\subset\mathfrak d\) because \(\mathfrak h\) is
locally nilpotent, which proves existence.

\smallskip
\noindent
\textup{d)}
This is a special case of \SGARefLink{sga3:xii:7.9}{XII~7.9}\textup{\SGARefLink{sga3:xii:7.9:item:d:02}{(d)}}.

% END inputs/XIV/tex/components/21_expose_XIV_section5_theorem51_through_lemma53_en.tex

% BEGIN inputs/XIV/tex/components/22_expose_XIV_corollary54_through_corollary56_en.tex
% Authority: Expo14.pdf, local pp. 24--25 / combined pp. 858--859.
% This closes Section 5.
% Exact next cursor: Section 6 and Theorem 6.1,
% local p. 25 / combined p. 859.

\paragraph{\textbf{Corollary 5.4.}}
\SGARefTarget{sga3:xiv:5.4}
{\itshape Let \(G\) be a smooth connected algebraic group over an
algebraically closed field \(k\). Let \(H\) be a connected
algebraic subgroup such that \(\mathfrak h\) contains a Cartan
subalgebra of \(\mathfrak g\), or equivalently such that \(H\)
contains a subgroup of type \((\mathrm C)\) of \(G\). Then the
number of conjugates of \(H\) containing a regular element of
\(G(k)\) equals the number of conjugates of \(\mathfrak h\)
containing a regular element of \(\mathfrak g\).}

Indeed, let \(g\) be a regular element of \(G(k)\), and let \(C\)
be the unique Cartan subgroup of \(G\) containing \(g\)
\textup{(XIII~2)}. The conjugates of \(H\) that contain \(g\) are
exactly those that contain \(C\)
\textup{(\SGARefLink{sga3:xiii:corollary:2:8}{XIII~2.8}\textup{(b)})}.

Similarly, let \(x\) be a regular element of \(\mathfrak g\). If a
conjugate \(\mathfrak h'\) of \(\mathfrak h\) contains \(x\), then
\[
  \operatorname{ad}(x)_{\mathfrak g/\mathfrak h'}
\]
is injective by \SGARefLink{sga3:xiii:corollary:5:4}{XIII~5.4}, and therefore
\[
  \mathfrak h'\supset\Nil(x,\mathfrak g)=\mathfrak d.
\]
Thus the number of conjugates of \(\mathfrak h\) containing \(x\)
equals the number containing the Cartan subalgebra
\(\mathfrak d\).

The subalgebra \(\mathfrak d\) is the Lie algebra of a subgroup
\(D\) of type \((\mathrm C)\) of \(G\). We may suppose that
\(C\subset D\). It remains to show that
\[
  H\supset C\quad\Longleftrightarrow\quad
  \mathfrak h\supset\mathfrak d.
\]
By \SGARefLink{sga3:xiii:corollary:5:5}{XIII~5.5}, the inclusion
\(\mathfrak h\supset\mathfrak d\) implies \(H\supset D\), and
\emph{a fortiori} \(H\supset C\). Conversely, by \SGARefLink{sga3:xiv:5.1}{5.1} the Lie
algebra \(\mathfrak t\) of \(C\) contains a regular element \(x\)
of \(\mathfrak g\). Hence \(H\supset C\) implies
\(x\in\mathfrak h\), and, as observed above, this implies
\(\mathfrak h\supset\mathfrak d\). This completes the proof.

Let \(G\) be as in \SGARefLink{sga3:xiv:5.2}{5.2}, and suppose that the infinitesimal rank of
its fibers is locally constant, that is, condition
\((\mathrm C_0)\) holds. By \SGARefLink{sga3:xiv:5.2}{5.2}\textup{(c)}, there is a morphism of
functors on \((\operatorname{Sch}/S)^\circ\),
\[
  \mathcal C\longrightarrow\mathcal D,
\]
where
\[
  \begin{aligned}
  \mathcal C(S')&=
    \{\text{Cartan subgroups of }G_{S'}\},\\
  \mathcal D(S')&=
    \{\text{subgroups of type \((\mathrm C)\) of }G_{S'}\}.
  \end{aligned}
\]
By \SGARefLink{sga3:xiv:3.10}{3.10} and \SGARefLink{sga3:xiv:5.2}{5.2}\textup{(a)}, the functor \(\mathcal D\) is
represented by a smooth quasi-projective prescheme over \(S\).
Let \(D\) be the universal subgroup of type \((\mathrm C)\) of
\(G_{\mathcal D}\). One may then define the functor
\[
  \mathcal C_D\colon
  (\operatorname{Sch}/\mathcal D)^\circ
  \longrightarrow\operatorname{Ens}
\]
for the \(\mathcal D\)-group \(D\), exactly as \(\mathcal C\) was
defined for the \(S\)-group \(G\).

\paragraph{\textbf{Proposition 5.5.}}
\SGARefTarget{sga3:xiv:5.5}
{\itshape Under the preceding hypotheses, regard \(\mathcal C\) as
a functor over \(\mathcal D\). Then \(\mathcal C\) is
\(\mathcal D\)-isomorphic to the functor \(\mathcal C_D\) of
Cartan subgroups of \(D\).}

This follows immediately from the next corollary.

\paragraph{\textbf{Corollary 5.6.}}
\SGARefTarget{sga3:xiv:5.6}
{\itshape Let \(G\) be a smooth \(S\)-prescheme in groups of finite
presentation, and let \(D\) be a subgroup of type \((\mathrm C)\)
of \(G\). There is a bijective correspondence between the Cartan
subgroups of \(G\) contained in \(D\) and the Cartan subgroups of
\(D\). More precisely, for a subgroup \(H\subset D\), \(H\) is a
Cartan subgroup of \(G\) if and only if it is a Cartan subgroup of
\(D\).}

This is a special case of \SGARefLink{sga3:xii:7.9}{XII~7.9}\textup{\SGARefLink{sga3:xii:7.9:item:c:02}{(c)}}, taking into account
that, over an algebraically closed field, every subgroup of type
\((\mathrm C)\) of \(G\) contains a Cartan subgroup of \(G\).

For the next section, the principal result obtained here is
\SGARefLink{sga3:xiv:5.2}{5.2}\textup{(c)} with \(H\) a Cartan subgroup of \(G\). It permits
the formulation of \SGARefLink{sga3:xiv:5.5}{5.5} and thus gives a useful dévissage of
\(\mathcal C\).

% END inputs/XIV/tex/components/22_expose_XIV_corollary54_through_corollary56_en.tex

% BEGIN inputs/XIV/tex/components/23_expose_XIV_section6_theorem61_en.tex
% Authority: Expo14.pdf, local pp. 25--27 / combined pp. 859--861.
% This begins Section 6 and contains Theorem 6.1.
% Exact next cursor: Corollary 6.2, local p. 27 / combined p. 861.

\section{Applications to the Structure of Algebraic Groups}
\SGARefTarget{sga3:xiv:6}

\paragraph{\textbf{Theorem 6.1.}}
\SGARefTarget{sga3:xiv:6.1}
{\itshape Let \(G\) be a smooth algebraic group over a field
\(k\). Consider the scheme \(\mathcal T\) of maximal tori of
\(G\), isomorphic to the scheme \(\mathcal C\) of Cartan
subgroups of \(G\). This is a homogeneous space under \(G^0\) and
a smooth affine connected algebraic scheme
\textup{(\SGARefLink{sga3:xii:7.1}{XII~7.1}\textup{\SGARefLink{sga3:xii:7.1:item:d}{(d)}})}. Then \(\mathcal C\) is a rational
variety: its field of rational functions is a purely
transcendental extension of \(k\).}

We first prove the theorem when \(k\) is infinite. We may plainly
suppose that \(G\) is connected, since \(\mathcal T\), and hence
\(\mathcal C\), is unchanged when \(G\) is replaced by \(G^0\).
Moreover, by \SGARefLink{sga3:xii:7.6}{XII~7.6}, \(\mathcal C\) is unchanged when \(G\) is
divided by a central subgroup. We may therefore first divide by
the center of \(G\) and suppose that \(G\) is affine
\textup{(\SGARefLink{sga3:xii:6.1}{XII~6.1})}, and then divide by its reductive center
\textup{(\SGARefLink{sga3:xii:4.1}{XII~4.1} and \SGARefLink{sga3:xiv:4.4}{4.4})} and suppose that the reductive center
of \(G\) is trivial \textup{(\SGARefLink{sga3:xii:4.7}{XII~4.7}\textup{\SGARefLink{sga3:xii:4.7:item:b}{(b)}})}.

We proceed by induction on \(n=\dim G\), assuming the result for
dimensions \(n'<n\). If \(G\) is nilpotent, then \(\mathcal C\)
consists of a single \(k\)-rational point and the assertion is
trivial. Otherwise, the Lie algebra of \(G\) is not nilpotent
\textup{(\SGARefLink{sga3:xiv:1.3}{1.3})}. Its Cartan subalgebras, and therefore the
subgroups of type \((\mathrm C)\) of \(G\), have some dimension
\(n'<n\). Consider the morphism
\[
  \mathcal C\longrightarrow\mathcal D
\]
of \SGARefLink{sga3:xiv:5.5}{5.5}. By \SGARefLink{sga3:xiv:3.10}{3.10}, since the infinite-field Lie algebra
\(\mathfrak g\) contains a regular element, \(\mathcal D\) is a
rational variety. Thus its function field \(K\) is purely
transcendental over \(k\).

Let \(x\) be the generic point of \(\mathcal D\). By \SGARefLink{sga3:xiv:5.5}{5.5}, the
fiber of \(\mathcal C\) at \(x\) is the scheme of Cartan
subgroups of a smooth connected algebraic group \(D_x\) over
\(K=\kappa(x)\), namely the generic subgroup of type
\((\mathrm C)\) of \(G\). The function field \(L\) of
\(\mathcal C\) is therefore the function field of
\(\mathcal C_{D_x}\). Since \(\dim D_x=n'<n\), the induction
hypothesis makes \(L\) purely transcendental over \(K\), and
hence, by transitivity, purely transcendental over \(k\).

Suppose now that \(k\) is finite. A different argument is
required. We may again suppose that \(G\) is affine and
connected. Since \(k\) is perfect, the radical of
\(G_{\bar k}\) descends to a subgroup \(R\) of \(G\).
Suppose first that \(R\ne G\), so that \(G\) is not solvable, and
let
\[
  u\colon G\longrightarrow G'=G/R
\]
be the canonical morphism. By \SGARefLink{sga3:xii:7.1}{XII~7.1}\textup{\SGARefLink{sga3:xii:7.1:item:e}{(e)}}, it induces
\[
  v\colon\mathcal C_G\longrightarrow\mathcal C_{G'},
  \qquad C\longmapsto u(C).
\]
Let \(x\) be the generic point of \(\mathcal C_{G'}\), and put
\(K=\kappa(x)\). The fiber \(v^{-1}(x)\) is the scheme of Cartan
subgroups of \(G_K\) whose image in \(G'_K\) is the generic
Cartan subgroup \(C'_x\) of \(G'\). By \SGARefLink{sga3:xii:7.9}{XII~7.9}\textup{\SGARefLink{sga3:xii:7.9:item:c:02}{(c)}}, it is
also the scheme of Cartan subgroups of
\[
  H=u_K^{-1}(C'_x).
\]
The field \(K\) is infinite. The infinite-field case just proved
therefore shows that the function field \(L\) of
\(\mathcal C_G\), which is also that of \(\mathcal C_H\), is
purely transcendental over \(K\).\footnote{The French source
refers here to the ``already proved part of \SGARefLink{sga3:xiv:6.1}{7.1}''; the argument
is the infinite-field part of \SGARefLink{sga3:xiv:6.1}{Theorem 6.1}.}

It remains to prove that \(K\) is purely transcendental over
\(k\), and we are thus reduced to the case in which \(G\) is
semisimple. Dividing by its reductive center, we may further
suppose that \(G\) is adjoint. Then \SGARefLink{sga3:xiv:3.18}{3.18} gives
\(\mathcal C_G\simeq\mathcal D_G\), and by \SGARefLink{sga3:xiv:3.10}{3.10} it is enough to
show that \(\mathfrak g\) has a regular point. As noted in \SGARefLink{sga3:xiv:3.20}{3.20},
this is an unpublished result of Chevalley.\footnote{See the
Appendix by J.-P. Serre.}

It remains only to treat the case in which \(k\) is finite and
\(G\) is connected, affine, and solvable. In fact, the next
result holds over an arbitrary field.

% END inputs/XIV/tex/components/23_expose_XIV_section6_theorem61_en.tex

% BEGIN inputs/XIV/tex/components/24_expose_XIV_corollary62_lemma63_en.tex
% Authority: Expo14.pdf, local pp. 27--28 / combined pp. 861--862.
% This completes the proof of Theorem 6.1.
% Exact next cursor: Corollary 6.4, local p. 28 / combined p. 862.

\paragraph{\textbf{Corollary 6.2.}}
\SGARefTarget{sga3:xiv:6.2}
{\itshape Let \(G\) be a smooth solvable algebraic group over a
field \(k\). Then the variety \(\mathcal C\) of Cartan subgroups
of \(G\) is isomorphic to an affine space
\(\operatorname{Spec}k[t_1,\ldots,t_N]\).}

We may suppose that \(G\) is connected and affine. Let
\(G_\infty\) be the terminal group in the descending central
series
\[
  G_0=G,\qquad G_{i+1}=[G,G_i].
\]
Thus \(G_\infty\) is the smallest invariant algebraic subgroup
of \(G\) for which \(G/G_\infty\) is nilpotent. Let \(C\) be a
Cartan subgroup of \(G\), which exists by \SGARefLink{sga3:xiv:1.1}{1.1}. Its image in
\(G/G_\infty\) is a Cartan subgroup and hence is all of
\(G/G_\infty\). Consequently the morphism
\[
  G_\infty\longrightarrow\mathcal C,\qquad
  g\longmapsto\operatorname{Ad}(g)C
\]
is an epimorphism and identifies
\[
  \mathcal C\simeq G_\infty/(N\cap G_\infty),
\]
where \(N\) is the normalizer of \(C\) in \(G\). In fact \(N=C\),
as recalled in \SGARefLink{sga3:xiv:4.4}{4.4}, but that is immaterial here.

The group \(U=G_\infty\) is smooth, connected, and unipotent:
over \(\bar k\) it lies in the unipotent part of \(G\), by the
known structure of smooth affine solvable groups
\textup{\emph{Bible}~6, Theorem~3}. If \(k\) is perfect, as in
the finite-field case needed for \SGARefLink{sga3:xiv:6.1}{6.1}, it follows easily that
\(U\) is \(k\)-unipotent: it has a composition series by
algebraic subgroups \(U_i\) with
\(U_i/U_{i+1}\simeq\mathbb G_a\).

Rosenlicht proved that a group of the form \(G_\infty\) above is
\(k\)-unipotent without any restriction on \(k\)
\textup{(M.~Rosenlicht, \emph{Questions of rationality for
solvable algebraic groups over non perfect fields}, Annali di
Matematica, 1963, pp.~97--120, \SGARefLink{sga3:xiv:7:lemma4}{Theorem~4}, \SGARefLink{sga3:xiv:7:lemma2}{Corollary~2})}. We admit
this substantially more delicate result and apply the following
lemma.

\paragraph{\textbf{Lemma 6.3.}}
\SGARefTarget{sga3:xiv:6.3}
{\itshape Let \(U\) be a smooth connected algebraic group over a
field \(k\), and let \(X=U/V\) be a homogeneous space under \(U\)
having a \(k\)-rational point. If \(U\) is \(k\)-unipotent, then,
as a \(k\)-scheme, \(X\) is isomorphic to an affine space
\(\operatorname{Spec}k[t_1,\ldots,t_N]\).}

Choose a composition series
\[
  U=U_0\supset U_1\supset\cdots\supset U_n=(e)
\]
by smooth connected normal subgroups of \(U\), with
\(U_i/U_{i+1}\simeq\mathbb G_a\). The groups
\[
  K_i=U_iV
\]
are algebraic subgroups of \(U\). They need not be smooth or
connected when \(V\) is not, but this does not matter, and
\(K_{i+1}\) is normal in \(K_i\). The canonical epimorphism
\[
  U_i/U_{i+1}\longrightarrow K_i/K_{i+1}
\]
shows that \(K_i/K_{i+1}\) is either the unit group or isomorphic
to \(\mathbb G_a/H_i\), where \(H_i\) is a finite subgroup of
\(\mathbb G_a\). Rosenlicht's \SGARefLink{sga3:xiv:7:lemma2}{Theorem~2} in the cited paper gives
\[
  K_i/K_{i+1}\simeq\mathbb G_a;
\]
this last assertion is immediate when \(k\) is perfect.

Put \(X_i=U/K_i\). We prove inductively that every \(X_i\) is an
affine space. If \(K_i/K_{i+1}\) is trivial, then
\(X_i=X_{i+1}\). Otherwise \(X_{i+1}\) is a principal
\(\mathbb G_a\)-bundle over \(X_i\). Since \(X_i\) is affine,
the bundle is trivial, so
\[
  X_{i+1}\simeq X_i\times\mathbb G_a.
\]
The induction proves the lemma. It proves \SGARefLink{sga3:xiv:6.2}{6.2} and thereby
completes the proof of \SGARefLink{sga3:xiv:6.1}{6.1}.

% END inputs/XIV/tex/components/24_expose_XIV_corollary62_lemma63_en.tex

% BEGIN inputs/XIV/tex/components/25_expose_XIV_corollary64_through_corollary610_en.tex
% Authority: Expo14.pdf, local pp. 28--30 / combined pp. 862--864.
% Exact next cursor: Remarks 6.11, local p. 30 / combined p. 864.

\paragraph{\textbf{Corollary 6.4.}}
\SGARefTarget{sga3:xiv:6.4}
{\itshape Let \(G\) be a smooth algebraic group over an infinite
field \(k\). Then the set of \(k\)-rational points of
\(\mathcal C\), in the notation of \SGARefLink{sga3:xiv:6.1}{6.1}, is Zariski dense. The
union of the Cartan subgroups of \(G\) is dense in \(G\).}

The first assertion holds for every rational variety over an
infinite field and is, for us, the most important arithmetic
consequence of the rationality result. The second follows from
the first and the density theorem \SGARefLink{sga3:xiii:theorem:2:1}{XIII~2.1}.

\paragraph{\textbf{Corollary 6.5.}}
\SGARefTarget{sga3:xiv:6.5}
{\itshape Let \(G\) be a smooth connected algebraic group over
\(k\). Then the variety \(Z\) of regular semisimple points of
\(G\) \textup{(\SGARefLink{sga3:xiii:corollary:3:5}{XIII~3.5})} is unirational. In particular, if
\(k\) is infinite, its \(k\)-rational points are dense in \(Z\).}

Indeed, \(Z\) is an open subscheme of a torus over
\(\mathcal C\). Its function field \(L\) is therefore the
function field of a torus over the function field \(K\) of
\(\mathcal C\), namely the generic maximal torus of \(G\). Thus
\(L\) is unirational over \(K\) by \SGARefLink{sga3:xiii:lemma:3:4}{XIII~3.4}. Since \SGARefLink{sga3:xiv:6.1}{6.1} makes
\(K\) purely transcendental over \(k\), \(L\) is unirational
over \(k\).

\paragraph{\textbf{Corollary 6.6.}}
\SGARefTarget{sga3:xiv:6.6}
{\itshape Let \(G\) be a smooth connected algebraic group over
\(k\), and let \(H\) be the subgroup generated by the subvariety
\(Z\) of regular semisimple points. Equivalently
\textup{(\SGARefLink{sga3:xii:8.2}{XII~8.2})}, \(H\) is the smallest invariant algebraic
subgroup of \(G\) for which \(G/H\) has reductive rank zero; over
\(\bar k\), \(G/H\) is an extension of an abelian variety by a
smooth connected unipotent group. Then \(H\) is unirational.

In particular, suppose \(G=H\); equivalently
\textup{(\SGARefLink{sga3:xii:8.4}{XII~8.4})}, suppose \(G\) is affine and there is no
nontrivial homomorphism \(G_{\bar k}\to\mathbb G_a\). Then \(G\)
is unirational. If, in addition, \(k\) is infinite, \(G(k)\) is
dense in \(G\).}

This follows immediately from \SGARefLink{sga3:xiv:6.5}{6.5}. Indeed, if \(Z_i\) are smooth
connected unirational \(k\)-preschemes and
\(u_i\colon Z_i\to G\) are morphisms, then the algebraic subgroup
of \(G\) generated by the \(u_i\) is unirational
\textup{(VI\textsubscript{B}~\SGARefLink{prop:sga3-VIB-7.1}{7.1})}.

An interesting special case of either \SGARefLink{sga3:xiv:6.5}{6.5} or \SGARefLink{sga3:xiv:6.6}{6.6} is the
following.

\paragraph{\textbf{Corollary 6.7.}}
\SGARefTarget{sga3:xiv:6.7}
{\itshape Let \(G\) be a smooth connected affine algebraic group
of unipotent rank zero. Then \(G\) is unirational.}

The conclusion of \SGARefLink{sga3:xiv:6.4}{6.4} may be sharpened as follows.

\paragraph{\textbf{Corollary 6.8.}}
\SGARefTarget{sga3:xiv:6.8}
{\itshape Let \(G\) be a smooth connected algebraic group over an
infinite field \(k\), and let \(C\) be a Cartan subgroup of \(G\).
Then the union of the conjugates of \(C\) by regular semisimple
elements of \(G(k)\) is dense in \(G\).}

This follows at once from \SGARefLink{sga3:xiv:6.4}{6.4} and \SGARefLink{sga3:xiii:corollary:3:6}{XIII~3.6}, which says that the
morphism
\[
  \varphi\colon Z\times C\longrightarrow G,\qquad
  \varphi(t,c)=\operatorname{Ad}(t)c,
\]
is dominant. The same result also implies the next corollary,
without any assumption that \(k\) is infinite.

\paragraph{\textbf{Corollary 6.9.}}
\SGARefTarget{sga3:xiv:6.9}
{\itshape Let \(G\) be a smooth connected algebraic group over
\(k\), and let \(H\) be a smooth connected algebraic subgroup
having the same reductive rank and the same nilpotent rank as
\(G\). Equivalently, over \(\bar k\), the group \(H_{\bar k}\)
contains a Cartan subgroup of \(G_{\bar k}\). If \(H\) is
unirational, so is \(G\). If \(H(k)\) is dense in \(H\), then
\(G(k)\) is dense in \(G\).}

The morphism
\[
  \varphi\colon Z\times H\longrightarrow G,\qquad
  \varphi(t,h)=\operatorname{Ad}(t)h,
\]
is dominant. By \SGARefLink{sga3:xiv:6.5}{6.5}, \(Z\) is unirational, and by hypothesis so
is \(H\), hence so are \(Z\times H\) and \(G\). The density
assertion is proved in the same way.

We recover the following well-known result, due to Chevalley in
characteristic zero and to Rosenlicht in characteristic
\(p>0\).

\paragraph{\textbf{Corollary 6.10.}}
\SGARefTarget{sga3:xiv:6.10}
{\itshape Let \(G\) be a smooth connected affine algebraic group
over a perfect field \(k\). Then \(G\) is unirational. If \(k\)
is also infinite, \(G(k)\) is dense in \(G\).}

By \SGARefLink{sga3:xiv:1.1}{1.1}, \(G\) has a Cartan subgroup \(C\). By \SGARefLink{sga3:xiv:6.9}{6.9}, it is enough
to prove that \(C\) is unirational. Since \(k\) is perfect,
Galois descent from the algebraically closed case
\textup{\emph{Bible}~6, \SGARefLink{sga3:xiv:7:lemma2}{Theorem~2}} gives
\[
  C=T\times C_u,
\]
where \(T\) is the maximal torus of \(C\) and \(C_u\) is smooth,
connected, and unipotent. The torus \(T\) is unirational by
XIII~3.4. Moreover, \(C_u\) is \(k\)-unipotent because \(k\) is
perfect, so \SGARefLink{sga3:xiv:6.3}{6.3} applies to \(C_u\).

% END inputs/XIV/tex/components/25_expose_XIV_corollary64_through_corollary610_en.tex

% BEGIN inputs/XIV/tex/components/26_expose_XIV_remarks611_en.tex
% Authority: Expo14.pdf, local pp. 30--31 / combined pp. 864--865.
% This closes Section 6, including the later editorial proof in note 8.
% Exact next cursor: Appendix, local p. 32 / combined p. 866.

\paragraph{\textbf{Remarks 6.11.}}
\SGARefTarget{sga3:xiv:6.11}
\textup{(a)} Rosenlicht gave examples of twisted forms of
\(\mathbb G_a\) over a nonperfect field having only finitely many
rational points; \emph{a fortiori}, they are not unirational.
Chevalley, on the other hand, gave an example of a torus over a
field of characteristic zero that is not rational. By contrast,
Chevalley's theory of semisimple groups shows that, over an
algebraically closed field, every smooth connected affine
algebraic group is rational. In any case, the question of
unirationality arises only for affine algebraic groups: by
Chevalley's structure theorem, a unirational algebraic group is
necessarily affine.

\textup{(b)} With the notation of \SGARefLink{sga3:xiv:6.6}{6.6}, it is tempting to seek a
criterion for the unirationality of \(G\) in terms of \(G/H\),
which is unipotent when \(G\) is affine. Clearly \(G/H\) must be
unirational; is that condition also sufficient? An example of
Rosenlicht shows that a smooth connected unipotent algebraic
group \(U\) may be rational without being \(k\)-unipotent.

\textup{(c)} Over a finite field \(k\), it would be interesting
to study questions of ``density'' such as the following one,
raised by Rosenlicht: if \(G\) is a smooth connected algebraic
group over \(k\), is \(G\) generated by its Cartan
subgroups?\footnote{This question was subsequently answered in
the affirmative by Steinberg.}

\begin{quote}\small
\textbf{Editors' note.}
The editors did not
identify this result in Steinberg's \emph{Collected Papers}; here
is a proof based on the Bruhat decomposition.

Let \(G/k\) be semisimple over the finite field \(k\), and let
\(G^\sharp\) be the subgroup generated by its \(k\)-subtori. The
question is stable under central isogeny and restriction of
scalars, so, as in \SGARefLink{sga3:xiv:7:lemma1}{Lemma~1} of the Appendix below, it is enough to
treat the geometrically simple case. The group \(G\) is
quasi-split by Lang's theorem and admits a Killing pair
\((T,B)\). By definition \(T\subset G^\sharp\), hence \(T\)
normalizes \(G^\sharp\). The big cell
\[
  R_u(B^-)\,T\,R_u(B)
\]
shows that it suffices to prove \(R_u(B)\subset G^\sharp\).
Let \(S\) be the maximal split \(k\)-torus of \(T\). For the
relative root system \(\Phi(G,S)\), choose a basis
\(\Delta_k\), and denote by \(U(\alpha)\) the unipotent subgroup
associated with \(\alpha\)
\textup{(A.~Borel, \emph{Linear Algebraic Groups}, second
edition, 1991, Proposition~21.9)}. Since \(R_u(B)\) is generated
by the \(U(\alpha)\), \(\alpha\in\Delta_k\), it is enough to
show \(U(\alpha)\subset G^\sharp\).

The classification supplies a quasi-split semisimple subgroup
\(G_\alpha\) of type \(A_1\), \({}^{2}A_1\), \({}^{3}A_1\), or
\({}^{2}A_2\), with
\[
  U(\alpha)\subset G_\alpha\subset G.
\]
Thus it is enough to consider \(G=\operatorname{PGL}_2\) or
\(G=\operatorname{SU}_3(K)\), where \(K/k\) is the unique
quadratic extension.

For \(\operatorname{PGL}_2\), since \(G^\sharp\) contains the
standard split torus \(T\), the possibilities up to
\(G(k)\)-conjugacy are
\[
  G^\sharp=T,\qquad G^\sharp=B,\qquad
  G^\sharp=\operatorname{PGL}_2.
\]
The first is excluded because \(G^\sharp\) contains the
\(k\)-torus \(R_{K/k}(\mathbb G_m)/\mathbb G_m\). The preceding
discussion also shows that \(G^\sharp=B\) forces
\(G^\sharp=\operatorname{PGL}_2\). Hence
\(G^\sharp=\operatorname{PGL}_2\).

For \(G=\operatorname{SU}_3(K)\), one has
\(\Delta_k=\{\alpha\}\), and the possibilities up to conjugacy
are
\[
 \begin{split}
  G^\sharp={}&T,\quad B,\quad U(2\alpha)\rtimes T,\\
             &\langle U(2\alpha),U(-2\alpha),T\rangle
               =\operatorname{SL}_2\cdot T,\quad\text{or }G.
 \end{split}
\]
The case \(B\) is excluded as for \(\operatorname{PGL}_2\). The
case \(\operatorname{SL}_2\cdot T\) is excluded because
\(G^\sharp\) contains the \(k\)-torus
\[
  R^1_{K/k}\!\left(R^1_{K_3/K}(\mathbb G_m)\right),
\]
where \(K_3/K\) is the degree-three extension, and the isogeny
class of this torus is irreducible. The same torus excludes
\(T\) and \(U(2\alpha)\rtimes T\). Hence \(G^\sharp=G\).
\end{quote}

For Rosenlicht's question one may reduce to the affine case by
dividing by the center. In the semisimple case, an affirmative
answer would also follow from a strengthening of Chevalley's
existence result for regular points mentioned in \SGARefLink{sga3:xiv:3.20}{3.20}: one would
need a regular element of \(\mathfrak g\) outside
\(\mathfrak h=\operatorname{Lie}(H)\), where \(H\ne G\) is a
prescribed smooth algebraic subgroup of an adjoint semisimple
group \(G\) over the finite field \(k\).

% END inputs/XIV/tex/components/26_expose_XIV_remarks611_en.tex

% BEGIN inputs/XIV/tex/components/27_expose_XIV_appendix_opening_lemma1_en.tex
% Authority: Expo14.pdf, local pp. 31--32 / combined pp. 865--866.
% This opens Serre's Appendix and runs through the setup before Lemma 2.
% Exact next cursor: Lemma 2, local p. 32 / combined p. 866.

\section{Appendix: Existence of Regular Elements over Finite Fields}
\SGARefTarget{sga3:xiv:7}

\begin{center}
\itshape by J.-P. Serre
\end{center}

Throughout the Appendix, \(k\) denotes a finite field,
\(\bar k\) its algebraic closure, and
\(\mathcal G=\operatorname{Gal}(\bar k/k)\). If
\(q=\operatorname{Card}(k)\), then \(\mathcal G\) is
topologically generated by the Frobenius element
\[
  F\colon x\longmapsto x^q.
\]
We shall prove the following theorem.\footnote{The theorem is due
to Chevalley. The editor wishes to express his gratitude to
American Express which, by misplacing a trunk of Chevalley's
manuscripts, obliged him to reconstruct the proof.}

\paragraph{\textbf{Theorem.}}
\SGARefTarget{sga3:xiv:7:theorem}
{\itshape Let \(G\) be an adjoint semisimple group defined over
\(k\), and let \(\mathfrak g\) be its Lie algebra. Then the
\(k\)-Lie algebra \(\mathfrak g(k)\) contains a regular element.}

\paragraph{\textbf{Remarks.}}
\textup{\SGARefLink{sga3:xiv:7:lemma1}{(1)}} Recall that ``adjoint'' means that the center of
\(G\), as a subgroup scheme, is trivial. In view of the Chevalley
seminar, this also means that, for a maximal torus \(T\) of
\(G\), the character group \(X(T)\) of \(T_{\bar k}\) is
generated by the set \(R\) of roots. In particular, the rank of
\(\mathfrak g\) equals \(\dim T\), the rank of \(G\).

\textup{\SGARefLink{sga3:xiv:7:lemma2}{(2)}} The editor does not know whether the adjoint
hypothesis is indispensable.

\paragraph{\textbf{Lemma 1.}}
\SGARefTarget{sga3:xiv:7:lemma1}
{\itshape It is enough to prove the theorem when \(G\) is
geometrically simple.}

Here geometrically simple means that \(G_{\bar k}\) has no smooth
normal subgroup scheme other than \(G_{\bar k}\) and \((e)\);
equivalently, the associated root system is irreducible.

We may first suppose that \(G\) is indecomposable over \(k\).
Then \(G_{\bar k}\) is a product of geometrically simple
components \(S\), permuted transitively by \(\mathcal G\). If
\(\mathcal H\subset\mathcal G\) fixes one component \(S\), and
\(K\) is its fixed field, then \(S\) is defined over \(K\), and
the standard restriction-of-scalars argument gives
\[
  G=R_{K/k}(S).
\]
Similarly,
\[
  \mathfrak g=R_{K/k}(\mathfrak s),
\]
where \(\mathfrak s\) is the Lie algebra of \(S\). If the theorem
holds for \(S\), there is an element
\[
  x\in\mathfrak s(K)=\mathfrak g(k)
\]
regular in \(\mathfrak s\), and one checks immediately that it
is regular in \(\mathfrak g\). This proves the lemma.

Henceforth assume that \(G\) is geometrically simple, and choose
a maximal torus \(T\) of \(G\). Let \(X\) be the character group
of \(T_{\bar k}\), let \(R\) be its root system, \(W\) its Weyl
group, and \(E\) the group of automorphisms of \(X\) preserving
\(R\). Then \(W\) is normal in \(E\).

For another maximal torus \(T'\), write \(X',R',W',E'\) for the
corresponding objects. If
\(y\in G(\bar k)\) satisfies
\[
  yT_{\bar k}y^{-1}=T'_{\bar k},
\]
then \(\operatorname{int}(y)\) identifies
\((X',R',W',E')\) with \((X,R,W,E)\). Replacing \(y\) changes
this identification by an element of \(W\).

Frobenius acts on \(X'\) while preserving \(R'\), and thus
defines an element \(f_{T'}\in E'\). Put \(f=f_T\). After the
above identification \(E'\simeq E\), the element \(f_{T'}\)
becomes an element \(f'\in E\), well defined up to
\(W\)-conjugacy. We now compare \(f'\) with \(f\).

% END inputs/XIV/tex/components/27_expose_XIV_appendix_opening_lemma1_en.tex

% BEGIN inputs/XIV/tex/components/28_expose_XIV_appendix_lemmas2_4_theorem_end_en.tex
% Authority: Expo14.pdf, local pp. 32--33 / combined pp. 866--867.
% Exact next cursor: Proof of Lemma 3, local p. 33 / combined p. 867.

\paragraph{\textbf{Lemma 2.}}
\SGARefTarget{sga3:xiv:7:lemma2}
{\itshape One has \(f'\equiv f\pmod W\). Conversely, every
\(f'\in E\) satisfying this condition arises from a maximal torus
\(T'\) of \(G\).}

Put \(\bar T=T_{\bar k}\) and \(\bar T'=T'_{\bar k}\), and choose
\(y\in G(\bar k)\) such that
\[
  y\bar T y^{-1}=\bar T'.
\]
Since \(y^q\bar T y^{-q}=\bar T'\), the element
\[
  n=y^{-1}y^q
\]
belongs to the normalizer \(N(T)(\bar k)\). The action of \(f'\)
on the points of \(\bar T\) is
\[
 \begin{aligned}
  f'(t)
    &=y^{-1}(yty^{-1})^q y\\
    &=y^{-1}y^q\,t^q\,y^{-q}y
      =nt^qn^{-1}.
 \end{aligned}
\]
If \(w\in W\) is represented by \(n\), then \(f'\) and
\(f\circ w\) have the same action on \(\bar T(\bar k)\), hence
on its characters. Thus \(f'=f\circ w\), and
\(f'\equiv f\pmod W\).

Conversely, represent a given \(w\in W\) by
\(n\in N(T)(\bar k)\). Lang's theorem gives
\[
  n=y^{-1}y^q
\]
for some \(y\in G(\bar k)\). Then
\(y\bar T y^{-1}\) is defined over \(k\), and the preceding
calculation shows that its corresponding element is
\(f\circ w\).

\paragraph{\textbf{Lemma 3.}}
\SGARefTarget{sga3:xiv:7:lemma3}
{\itshape Let \(X,R,W,E\) be as above, with \(R\) irreducible,
and let \(\varphi\in E/W\). There exist a representative
\(f\in E\) of \(\varphi\) and roots
\(\theta_1,\ldots,\theta_n\) such that:
\begin{enumerate}
\item \((\theta_1,\ldots,\theta_n)\) is a basis of \(X\);
\item \(R\) is the union of the orbits of the \(\theta_i\) under
the powers of \(f\): every \(a\in R\) has the form
\(a=f^m\theta_i\) for suitable \(m\) and \(i\).
\end{enumerate}}

The proof will be given below. We first need a lemma of linear
algebra.

\paragraph{\textbf{Lemma 4.}}
\SGARefTarget{sga3:xiv:7:lemma4}
{\itshape Let \(V\) be a vector space over \(k\), and let
\((\theta_1,\ldots,\theta_n)\) be a basis of the dual of
\(V\otimes_k\bar k\). Then there is an \(x\in V\) such that
\(\theta_i(x)\ne0\) for every \(i\).}

Let \(V^*\) be the dual of \(V\), and let \(W_i\) be the
\(\bar k\)-subspace of \(V^*\otimes_k\bar k\) generated by the
Galois conjugates of \(\theta_i\). It descends to a
\(k\)-subspace, again denoted \(W_i\). The natural map
\[
  W_1\otimes\cdots\otimes W_n
     \longrightarrow\bigwedge\nolimits^n V^*
\]
is nonzero: after extension to \(\bar k\), it contains
\(\theta_1\wedge\cdots\wedge\theta_n\ne0\). Hence \(V^*\) has a
basis \(u_1,\ldots,u_n\) with \(u_i\in W_i\). Choose
\(x\in V\) with \(u_i(x)\ne0\) for all \(i\), for example with
\(u_i(x)=1\). If \(\theta_i(x)=0\), every Galois conjugate of
\(\theta_i\) vanishes at \(x\), and so does \(u_i\), a
contradiction.

\paragraph{\textbf{End of the proof of the theorem.}}

By \SGARefLink{sga3:xiv:7:lemma2}{Lemmas 2} and 3, choose a maximal torus \(T\) whose element
\(f\) has the properties of \SGARefLink{sga3:xiv:7:lemma3}{Lemma 3}. Let \(Y\) be the dual of
\(X\). The Lie algebra of \(T\) satisfies
\[
  \mathfrak t(\bar k)\simeq Y\otimes_{\mathbb Z}\bar k,
\]
compatibly with the Galois action. Put
\(V=\mathfrak t(k)\). An element \(x\in V\) is regular precisely
when it is annihilated by no root \(\alpha\in R\), viewed as a
linear form in \(V^*\otimes_k\bar k\).

\SGARefLink{sga3:xiv:7:lemma4}{Lemma 4} supplies \(x\in V\) annihilated by none of the
\(\theta_i\). Every root is a Galois conjugate of one of the
\(\theta_i\), by property \textup{\SGARefLink{sga3:xiv:7:lemma2}{(2)}} of \SGARefLink{sga3:xiv:7:lemma3}{Lemma 3}. Hence no root
annihilates \(x\), and \(x\) is regular. This proves the theorem.

% END inputs/XIV/tex/components/28_expose_XIV_appendix_lemmas2_4_theorem_end_en.tex

% BEGIN inputs/XIV/tex/components/29_expose_XIV_appendix_lemma3_proof_through_case2_en.tex
% Authority: Expo14.pdf, local pp. 33--34 / combined pp. 867--868.
% Exact next cursor: Case (3) in the proof of Lemma 3,
% local p. 34 / combined p. 868.

\paragraph{\textbf{Proof of \SGARefLink{sga3:xiv:7:lemma3}{Lemma 3}.}}

The proof uses Coxeter transformations. Let
\(a_1,\ldots,a_n\) be a simple system of roots of \(R\), and let
\(r_i\) be the reflection corresponding to \(a_i\). Put
\[
  c=r_1\cdots r_n.
\]
Then \(c\in W\). Although \(c\) depends on the simple system and
on its ordering, its conjugacy class does not. It is called the
Coxeter element. We shall use, without proof, the fact that \(1\)
is not an eigenvalue of \(c\).

\paragraph{\textbf{Lemma 5.}}
\SGARefTarget{sga3:xiv:7:lemma5}
Put
\[
  \theta_i=r_n\cdots r_{i+1}(a_i).
\]
Then:
\begin{enumerate}
\item[\textup{\SGAInlineRefTarget{sga3:xiv:7:lemma5:item:a}(a)}] the \(\theta_i\) form a basis of the group
\(X\) generated by the roots;
\item[\textup{\SGAInlineRefTarget{sga3:xiv:7:lemma5:item:b}(b)}] \(\theta_i>0\) and \(c(\theta_i)<0\) for
every \(i\);
\item[\textup{\SGAInlineRefTarget{sga3:xiv:7:lemma5:item:c}(c)}] every root \(a\in R\) such that \(a>0\) and
\(c(a)<0\) is one of the \(\theta_i\);
\item[\textup{\SGAInlineRefTarget{sga3:xiv:7:lemma5:item:d}(d)}] \(R\) is the union of the orbits of the
\(\theta_i\) under the powers of \(c\).
\end{enumerate}

Indeed,
\[
  \theta_i=a_i+\sum_{j>i}m_{ij}a_j,\qquad m_{ij}\in\mathbb Z,
\]
which proves \textup{\SGARefLink{sga3:xiv:7:lemma5:item:a}{(a)}}. Assertions \textup{\SGARefLink{sga3:xiv:7:lemma5:item:b}{(b)}} and
\textup{\SGARefLink{sga3:xiv:7:lemma5:item:c}{(c)}} follow from the observation that \(r_i\) preserves
the sign of every root other than \(\pm a_i\), and reverses the
sign of \(\pm a_i\).

For \textup{\SGARefLink{sga3:xiv:7:lemma5:item:d}{(d)}}, an orbit of \(c\) cannot consist entirely of
positive roots, nor entirely of negative roots. Otherwise the
sum of the roots in that orbit would give a nonzero element of
\(X\) fixed by \(c\), contrary to the fact that \(1\) is not an
eigenvalue. Hence every orbit contains a root \(a>0\) for which
\(c(a)<0\), and \textup{\SGARefLink{sga3:xiv:7:lemma5:item:c}{(c)}} applies.

\paragraph{\textbf{Remark.}}
The proof was recalled only for the reader's convenience; one
could instead refer to the standard texts on Coxeter
transformations, for example Koszul, Séminaire Bourbaki
1959/1960, Exposé~191. They also show that all orbits of \(c\)
have the same number \(h\) of elements and each contains exactly
one \(\theta_i\). In particular,
\[
  h=\frac{\operatorname{Card}(R)}{n}.
\]

We return to \SGARefLink{sga3:xiv:7:lemma3}{Lemma 3} and distinguish three cases.

\medskip
\noindent\textup{\SGARefLink{sga3:xiv:7:lemma1}{(1)}}
Suppose that \(\varphi\in E/W\) is the identity. We must choose
\(f\in W\); take \(f=c\). \SGARefLink{sga3:xiv:7:lemma5}{Lemma 5} shows that this works.

\medskip
\noindent\textup{\SGARefLink{sga3:xiv:7:lemma2}{(2)}}
Suppose that \(R\) has type \(A_n\), where \(n\ge2\) is even, and
that \(\varphi\) is the unique nontrivial element of \(E/W\).
Then \(-1\notin W\), and we take
\[
  f=-c.
\]
The Coxeter element \(c\) has order \(h=n+1\), which is odd. If
\(a\in R\), \SGARefLink{sga3:xiv:7:lemma5}{Lemma 5} gives \(a=c^m\theta_i\) for suitable \(m,i\).
Adding \(h\) to \(m\) if necessary, we may suppose that \(m\) is
even. Then
\[
  a=(-c)^m\theta_i=f^m\theta_i.
\]
Thus the orbits of the \(\theta_i\) under \(f\) fill \(R\).

% END inputs/XIV/tex/components/29_expose_XIV_appendix_lemma3_proof_through_case2_en.tex

% BEGIN inputs/XIV/tex/components/30_expose_XIV_appendix_lemma3_case3_end_en.tex
% Authority: Expo14.pdf, local pp. 34--35 / combined pp. 868--869.
% This closes Lemma 3 and the Appendix.
% Local p. 36 / combined p. 870 is terminal blank.
% Exact next cursor: END OF EXPOSÉ XIV; Exposé XV begins combined p. 871.

\medskip
\noindent\textup{\SGARefLink{sga3:xiv:7:lemma3}{(3)}}
Suppose that \(\varphi\in E/W\) is nontrivial and that \(R\) has
one of the following types:
\[
  A_n\ (n\ge3\text{ odd}),\qquad
  D_4,\qquad D_n\ (n\ge5),\qquad E_6.
\]
These, together with the even \(A_n\) case treated separately
above, are exactly the cases in which \(E/W\) is nontrivial. The
following diagrams record the corresponding Dynkin-diagram
automorphisms.

\begin{center}
\begin{tikzpicture}[
  x=1.05cm,
  y=.92cm,
  >=Latex,
  dynkin node/.style={circle,fill=black,inner sep=2.2pt},
  dynkin edge/.style={line width=.45pt},
  dynkin arrow/.style={line width=.45pt}
]
  % A_n, n odd: reflection of the chain.
  \node[anchor=east] at (-1.4,6.7) {\(A_n,\ n\text{ odd}\geq 3\)};
  \foreach \x/\name in {0/aone,1.3/atwo,2.6/athree,3.9/afour,5.2/afive}
    \node[dynkin node] (\name) at (\x,6.7) {};
  \draw[dynkin edge] (aone)--(atwo)--(athree)--(afour)--(afive);
  \draw[dynkin arrow,<->]
    (aone) .. controls +(1.0,1.35) and +(-1.0,1.35) .. (afive);
  \draw[dynkin arrow,<->]
    (atwo) .. controls +(.65,.72) and +(-.65,.72) .. (afour);

  % D_4: the three outer vertices are cyclically permuted.
  \node[anchor=east] at (-1.4,4.3) {\(D_4\)};
  \node[dynkin node] (dcenter) at (2.6,4.3) {};
  \node[dynkin node] (dtop)    at (1.3,5.15) {};
  \node[dynkin node] (dbottom) at (1.3,3.45) {};
  \node[dynkin node] (dright)  at (4.0,4.3) {};
  \draw[dynkin edge] (dcenter)--(dtop);
  \draw[dynkin edge] (dcenter)--(dbottom);
  \draw[dynkin edge] (dcenter)--(dright);
  \draw[dynkin arrow,->]
    (dright) .. controls +(0,1.05) and +(.75,.48) .. (dtop);
  \draw[dynkin arrow,->]
    (dtop) .. controls +(-1.15,-.18) and +(-1.15,.18) .. (dbottom);
  \draw[dynkin arrow,->]
    (dbottom) .. controls +(.75,-.48) and +(0,-1.05) .. (dright);

  % D_n, n >= 5: interchange of the two terminal vertices.
  \node[anchor=east] at (-1.4,1.85) {\(D_n,\ n\geq 5\)};
  \foreach \x/\name in {0/dnone,1.25/dntwo,2.5/dnthree,3.75/dnbranch}
    \node[dynkin node] (\name) at (\x,1.85) {};
  \node[dynkin node] (dnupper) at (4.85,2.65) {};
  \node[dynkin node] (dnlower) at (4.85,1.05) {};
  \draw[dynkin edge] (dnone)--(dntwo)--(dnthree)--(dnbranch);
  \draw[dynkin edge] (dnbranch)--(dnupper);
  \draw[dynkin edge] (dnbranch)--(dnlower);
  \draw[dynkin arrow,<->]
    (dnupper) .. controls +(.8,-.3) and +(.8,.3) .. (dnlower);

  % E_6: reflection across the central vertical axis.
  \node[anchor=east] at (-1.4,-.75) {\(E_6.\)};
  \foreach \x/\name in {0/eone,1.3/etwo,2.6/ethree,3.9/efour,5.2/efive}
    \node[dynkin node] (\name) at (\x,-.75) {};
  \node[dynkin node] (esix) at (2.6,-1.9) {};
  \draw[dynkin edge] (eone)--(etwo)--(ethree)--(efour)--(efive);
  \draw[dynkin edge] (ethree)--(esix);
  \draw[dynkin arrow,<->]
    (eone) .. controls +(1.0,1.35) and +(-1.0,1.35) .. (efive);
  \draw[dynkin arrow,<->]
    (etwo) .. controls +(.65,.72) and +(-.65,.72) .. (efour);
\end{tikzpicture}
\end{center}

Let \(S\) be a simple system of roots, not yet ordered. The group
\(E\) is the semidirect product of \(W\) with the group
\(\Psi\) of permutations of \(S\) preserving the Cartan matrix,
or equivalently the Dynkin diagram. Thus \(E/W\simeq\Psi\), and
\(\varphi\) corresponds to some \(\psi\in\Psi\).

Inspection of the diagrams shows that each \(\psi\)-orbit in
\(S\) consists of pairwise orthogonal roots. This fails for
\(A_n\) with \(n\) even, which is why that case was separated.
If \(\sigma\) is such an orbit, the reflections \(r_a\),
\(a\in\sigma\), commute; write
\[
  \rho_\sigma=\prod_{a\in\sigma}r_a.
\]
Plainly \(\rho_\sigma\) commutes with \(\psi\).

Choose a total order on \(S\) for which every \(\psi\)-orbit is a
segment, and write \(S=\{a_1,\ldots,a_n\}\). Let
\(\theta_1,\ldots,\theta_n\) and
\(c=r_1\cdots r_n\) be defined as above. The element \(c\) is a
product of the \(\rho_\sigma\), with the orbits in some order, so
it commutes with \(\psi\). Put
\[
  f=c\psi.
\]

The automorphism \(\psi\) permutes the \(\theta_i\). Indeed,
\(\psi(\theta_i)>0\), and, since \(c\) commutes with \(\psi\),
\[
  c(\psi(\theta_i))=\psi(c(\theta_i))<0.
\]
\SGARefLink{sga3:xiv:7:lemma5}{Lemma 5}\textup{\SGARefLink{sga3:xiv:7:lemma5:item:c}{(c)}} therefore gives
\(\psi(\theta_i)=\theta_j\) for some \(j\).

Finally, if \(a\in R\), \SGARefLink{sga3:xiv:7:lemma5}{Lemma 5} gives
\(a=c^m\theta_i\) for some \(m,i\). Hence
\[
  a=f^m(\psi^{-m}\theta_i)=f^m\theta_j
\]
for a suitable \(j\). Thus \(f\) and the \(\theta_i\) satisfy
both conditions of \SGARefLink{sga3:xiv:7:lemma3}{Lemma 3}, completing its proof.

\paragraph{\textbf{Remark.}}
Except in case \textup{\SGARefLink{sga3:xiv:7:lemma2}{(2)}}, one can prove that every orbit of
\(f\) has exactly
\[
  h=\frac{\operatorname{Card}(R)}{n}
\]
elements and contains exactly one \(\theta_i\). In case
\textup{(2)}, some of the \(\theta_i\) are superfluous.

% END inputs/XIV/tex/components/30_expose_XIV_appendix_lemma3_case3_end_en.tex


